\chapter{Associative chiral algebras}
\label{e1f:chapter}

The scalar vacuum has an associative field product: multiplying
any number of its identity fields gives the identity. Its collision
residues\phantomsection\label{e1f:def:scalar-three-point-coordinates} behave differently. For three positions, write
\(u=z_1-z_3\) and \(v=z_2-z_3\). The function
\[
 F(u,v)=\frac1{u(u-v)}
\]
has double residue \(1\) when the first two positions meet first,
and double residue \(0\) when the last two meet first. The third
order, in which the outer positions meet with the middle one held
apart, contributes \(-1\). The three terms satisfy a Jacobi
relation; the first two do not satisfy associativity.

This distinction determines which coefficient objects can retain
an associative structure. A chiral operation in the sense of Beilinson and
Drinfeld takes localized differential-module inputs to an output
on their common diagonal \cite{chhop:BD}. On a curve, a strictly
associative operation of this type has zero triple products.
A normalized local vacuum then forces the module itself to vanish.
Replacing ordinary maps by derived maps on the same degree-zero
modules does not remove this obstruction.

There is nevertheless an associative construction retaining the
fields and their vacuum. Put their residues in the differential
of a coalgebra on configurations, and take its associative cobar.
The coefficients of this algebra are complexes on collision
strata. At three positions they retain the outer collision,
which supplies the homotopy between the two consecutive residue
operations. The unit of the cobar lies on the empty configuration;
the vacuum remains a state at one position.

The obstruction on modules makes the change of coefficient
objects necessary. Ordered
matrix-valued fields then determine the configuration geometry,
and the scalar vacuum determines its first nontrivial collision
coefficients. These examples lead to the residue coalgebra and
its cobar, to the comparison with symmetric chiral algebras, and
to the extra tensor and duality data required for coproducts and
centres. Throughout, \(k\) has characteristic zero and \(X\) is a
smooth curve over \(k\). Differential modules are right modules;
\(\omega_X\) denotes the canonical right \(\mathcal D_X\)-module.
Complexes are cohomologically graded, with
\(H^i(M[r])=H^{i+r}(M)\), and tensor operations on
homogeneous coefficients use the Koszul signs.

\section{Operations with a prescribed collision target}
\label{e1f:sec:pch}

For\phantomsection\label{e1f:def:distinct-configuration-open} a nonempty finite set \(I\), let
\(j_I:U_I\hookrightarrow X^I\) be the locus of pairwise distinct
positions, and let \(\Delta_I:X\hookrightarrow X^I\) be the small
diagonal. Given right differential modules \(M_i\) and \(N\), put
\begin{equation}\label{e1f:eq:pch}
 P_I^{\mathrm{ch}}(\{M_i\};N)
 =\operatorname{Hom}_{\mathcal D_{X^I}}
 \left(j_{I*}j_I^*\!\left(\boxtimes_{i\in I}M_i\right),
                    \Delta_{I*}N\right).
\end{equation}
The source permits poles where positions meet. The target is
supported where all positions agree, and retains the normal
derivatives of that diagonal. Taking an operation therefore
specifies both its singular inputs and its collision output.

For a surjection \(\pi:I\twoheadrightarrow J\), first apply an
operation inside each fibre. Localize between different fibres,
transfer their outputs along the partial diagonal, and apply the
operation indexed by \(J\). These are the compositions of the
chiral pseudo-tensor category \cite{chhop:BD}*{\S3.1}. They use
all surjections. Permuting the labels permutes the coordinates
and the external factors, so the category is symmetric.
A \emph{chiral algebra} is a Lie algebra for these operations:
its binary map is skew under this simultaneous permutation and
satisfies Jacobi.

Write\phantomsection\label{e1f:def:binary-open-diagonal} \(j=j_{\{1,2\}}\) and \(\Delta=\Delta_{\{1,2\}}\).
The principal-parts quotient is the chiral operation
\begin{equation}\label{e1f:eq:unit}
 \varepsilon_N:
 j_*j^*(\omega_X\boxtimes N)\twoheadrightarrow\Delta_*N.
\end{equation}
A vacuum of a chiral algebra \(A\) is a map
\(u:\omega_X\to A\) whose left action is \(\varepsilon_A\).
This normalization fixes the residue, rather than an ordinary
product of two sections at the same point.

On\phantomsection\label{e1f:def:translation-state-space} \(X=\mathbb A^1\), a vacuum and a compatible translation
action identify chiral algebras with vertex algebras
\cite{chhop:BD}*{\S0.15}. The translation action makes the
underlying differential module weakly equivariant. Let \(V\)
be its state space and \(T\) its translation on states. In a
translation trivialization the left-module form is
\(\mathcal O_X\otimes_k V\), with connection
\[
 \partial_z(f\otimes a)=f'\otimes a+f\otimes Ta.
\]
The\phantomsection\label{e1f:def:state-field-expansion} associated fields have the expansion
\[
 Y(a,u)b=\sum_{n\in\mathbb Z}(a_{(n)}b)u^{-n-1}.
\]
The chiral map encodes this whole expansion
\cite{chhop:BD}*{\S3.5.10}. Its value on an input multiplied by
\(u^{-j-1}\), followed by the simple-residue projection, extracts
the coefficient of \(u^j\). Restricting the inputs to be regular
along the diagonal retains only the singular coefficients.

The vacuum does not supply the missing regular coefficients.
For example, put \(T=0\) on \(k1\oplus kx\). The two unital
commutative products \(x^2=0\) and \(x^2=x\) have zero singular
part and the same vacuum. Their constant fields, and hence their
chiral maps on localized inputs, differ.

\subsection{The same operations on Ran space}

The\phantomsection\label{e1f:def:ran-singleton-system} Ran category \(\mathcal D(\operatorname{Ran}X)\) consists of
systems of differential-module complexes on \(X^I\), compatible
with exceptional pullback along every surjection. Equip it with
the chiral tensor of Francis and Gaitsgory
\cite{chhop:FG}*{\S\S2.1--2.3}, and let
\(\iota:X\to\operatorname{Ran}X\) be the singleton inclusion.
Its direct image \(\iota_*\) is fully faithful.

\begin{proposition}\label{e1c:prop:ran}
For differential modules \(M_1,\ldots,M_n,N\) in degree zero,
the degree-zero maps
\[
 \iota_*M_1\otimes^{\mathrm{ch}}\cdots
       \otimes^{\mathrm{ch}}\iota_*M_n\longrightarrow\iota_*N
\]
are precisely \(P_n^{\mathrm{ch}}(M_1,\ldots,M_n;N)\).
The identification preserves permutations and compositions along
surjections.
\end{proposition}

\begin{proof}
The tensor of the singleton objects is induced from
\(j_{n*}j_n^*(\boxtimes_iM_i)\) on \(X^n\)
\cite{chhop:FG}*{\S2.3.3, Lemma~2.3.4}. Chart adjunction and
base change identify its maps to \(\iota_*N\) with maps to
\(\Delta_{n*}N\). Both coefficient objects are concentrated in
degree zero, so degree-zero derived maps are ordinary module maps.
This gives \eqref{e1f:eq:pch}.

For an iterated operation, localization first within the fibres
and then between them is the localization of the full input.
Base change along the partial diagonal identifies the resulting
map with the surjection composition. Permutations act on the
same coordinates and external factors in both descriptions.
\end{proof}

The proposition concerns singleton-supported objects. The free
nonunital associative Ran algebra on \(\iota_*\omega_X\), with
underlying object
\[
 \bigoplus_{n\geq1}(\iota_*\omega_X)^{\otimes^{\mathrm{ch}}n},
\]
has associative concatenation, but its second tensor power is
nonzero on two distinct positions. It is not a singleton
insertion.

Adjoining the empty configuration gives a tensor unit
\(k_\varnothing\) supported there. A map from this unit to an
object supported entirely on nonempty configurations is zero;
the unit identity would make a unital such object zero.
Unitalization therefore adds the summand \(k_\varnothing\).
Its unit is different from the local map \(\omega_X\to A\).

\section{Why a vacuum excludes strict associativity}
\label{e1f:sec:collapse}

Call a binary operation \(\mu\in P_2^{\mathrm{ch}}(A,A;A)\)
\emph{strictly associative} when
\(\mu\circ_1\mu=\mu\circ_2\mu\) in
\(P_3^{\mathrm{ch}}(A,A,A;A)\). A map \(u:\omega_X\to A\)
is normalized on the left if
\(\mu(u,1)=\varepsilon_A\), and on the right if
\(\mu(1,u)=-\varepsilon_A\tau\). Here \(\tau\) exchanges
the factors and coordinates. The minus sign is the orientation
sign for transposing the two positions.

Before imposing either normalization, consider the scalar test
from the beginning of the chapter.

\begin{proposition}\label{e1f:prop:flags}
For \(F(u,v)=1/(u(u-v))\), taking first the residue at
\(u-v=0\) and then at \(v=0\) gives \(1\). Taking first the
residue at \(v=0\) and then at \(u=0\) gives \(0\).
\end{proposition}

\begin{proof}
Set \(s=u-v\), \(t=v\) for the first flag. In the domain
where \(s\) is expanded first,
\[
 \frac1{s(s+t)}
   =\sum_{n\geq0}(-1)^n s^{n-1}t^{-n-1}.
\]
Its coefficient of \(s^{-1}t^{-1}\) is \(1\).
For the other flag, put \(s=v\), \(t=u\). The expansion
\(1/(t(t-s))=\sum_{n\geq0}s^nt^{-n-2}\) has no negative
power of \(s\), so its first residue is zero.
\end{proof}

The remaining flag begins with \(u=0\), the collision of
positions \(1,3\). Expanding
\(F=-v^{-1}u^{-1}(1-u/v)^{-1}\) gives successive residues
\(-v^{-1}\) and \(-1\). Thus the three values are
\(1,0,-1\). Their signed relation involves the intermediate
stratum \(13|2\). It does not yet give a degree-minus-one
operation: that requires a complex containing this stratum.

\subsection{Separating the two consecutive channels}

The first two flags cannot agree nontrivially even when their
intermediate coefficients differ. The reason is visible before
the last collision. Along \(12|3\), the first normal coordinate
acts locally nilpotently; along \(1|23\), the second does. The
outer difference is invertible on both intermediate modules.

\begin{lemma}\label{e1f:lem:separation}
Let \(R\) be commutative, \(s,v\in R\), and \(h=s+v\).
Let \(E,L_s,L_v\) be \(R[h^{-1}]\)-modules, with \(s\)
locally nilpotent on \(L_s\) and \(v\) locally nilpotent on
\(L_v\). Let \(Q\) be an \(R\)-module. Suppose the maps
\[
 E\xrightarrow{\alpha_s}L_s\xrightarrow{\beta_s}Q,
 \qquad
 E\xrightarrow{\alpha_v}L_v\xrightarrow{\beta_v}Q
\]
have the same composite, where the \(\alpha\)'s are
\(R[h^{-1}]\)-linear and the \(\beta\)'s are \(R\)-linear.
Then their common composite is zero.
\end{lemma}

\begin{proof}
Fix \(e\in E\) and choose \(a,b\geq1\) such that
\(s^a\alpha_s(e)=0\) and \(v^b\alpha_v(e)=0\). With
\(N=a+b-1\), expand the identity
\[
 e=h^{-N}(s+v)^Ne
   =\sum_{i=0}^N\binom Ni s^iv^{N-i}h^{-N}e.
\]
For \(i\geq a\), calculate the common image through \(L_s\).
It is zero because \(\alpha_s\) commutes with \(h^{-N}\).
For \(i<a\), one has \(N-i\geq b\), and calculation through
\(L_v\) gives zero. All summands therefore vanish.
The inverse of \(h\) was used only before the maps to \(Q\).
\end{proof}

The exponent cannot be reduced uniformly. For independent
\(s,v\) over \(k\), the polynomial \((s+v)^{a+b-2}\)
contains the nonzero term
\(\binom{a+b-2}{a-1}s^{a-1}v^{b-1}\) modulo
\((s^a,v^b)\). Particular specializations may need a smaller
exponent; the lemma needs no uniform nilpotence bound.

\begin{theorem}\label{e1f:thm:twostep}
In the ordinary chiral pseudo-tensor category of a smooth curve,
the linear spans of the ternary channel images through \(12|3\)
and \(1|23\) have zero intersection. The intermediate modules
and binary maps may differ. In particular, a binary chiral
operation is strictly associative exactly when
\[
 \mu\circ_1\mu=0=\mu\circ_2\mu.
\]
Every tree operation built from such a \(\mu\) on at least
three inputs vanishes. No vacuum or uniform pole bound is needed.
\end{theorem}

\begin{proof}
The common target is supported on the small diagonal, so work
there in an \'etale coordinate. Put
\(s=z_1-z_2\), \(v=z_2-z_3\), and \(h=z_1-z_3=s+v\).
The fully localized input is
\(E=j_{3*}j_3^*(M_1\boxtimes M_2\boxtimes M_3)\).
The first map in the channel \(12|3\) lands in the direct
image from that partial diagonal with the third position kept
separate. It is linear over \(\mathcal O[h^{-1}]\).
Every section of this intermediate direct image uses finitely
many normal derivatives, and hence is killed by a power of
\(s\). The last map is \(\mathcal O\)-linear.
The other channel has the same properties with \(v\) in place
of \(s\).

A finite sum of channel maps factors through the direct sum of
its intermediate modules. Such a direct sum retains the stated
local nilpotence and localization. Lemma~\ref{e1f:lem:separation}
therefore proves the intersection assertion. Associativity
puts its two composites in that intersection. Conversely, two
zero composites are equal. Finally, every binary tree with at
least three leaves contains adjacent internal vertices; their
composition is one of these zero ternary maps.
\end{proof}

\subsection{Inserting the vacuum}

The scalar flag calculation suggests that two vacuum insertions
should detect every state. The following decomposition of
operation spaces makes that argument independent of coordinates.

\begin{lemma}[\cite{chhop:BD}*{Lemma~3.1.13}]
\label{e1f:lem:insert}
Let \(E_i\phi\) denote insertion of the principal-parts map
into the \(i\)-th input of \(\phi\). For nonempty \(I\),
\[
 \bigoplus_{i\in I}P_I^{\mathrm{ch}}(\{M_i\};N)
 \xrightarrow{\ \sim\ }
 P_{\{0\}\sqcup I}^{\mathrm{ch}}(\omega_X,\{M_i\};N),
 \qquad (\phi_i)\longmapsto\sum_iE_i\phi_i.
\]
Moreover, \(\varepsilon_N(1,\phi)=\sum_iE_i\phi\).
\end{lemma}

\begin{proof}
Put \(N_I=j_{I*}j_I^*(\boxtimes_iM_i)\), and let
\(\Delta^i\) be the graph \(x_0=x_i\). Over \(U_I\) these
graphs are disjoint. Taking principal parts in \(x_0\) gives
\[
 0\longrightarrow\omega_X\boxtimes N_I
 \longrightarrow j_*j^*(\omega_X\boxtimes\boxtimes_iM_i)
 \longrightarrow\bigoplus_i\Delta^i_*N_I
 \longrightarrow0.
\]
A map from the first term to \(\Delta_*N\) is zero. Indeed,
its generators \(dx_0\boxtimes n\) are killed by right
\(\partial_{x_0}\), whereas that operator is injective on
the target: it raises a normal derivative degree.
Every map from the middle term therefore factors uniquely
through the quotient. Transfer along the closed graphs
identifies its summands with the operation spaces on the left.
The principal part of \((x_0-x_i)^{-1}\) occurs only on the
\(i\)-th graph. Adding these graph contributions proves the
last identity.
\end{proof}

\begin{theorem}\label{e1f:thm:vacuum}
Let \(\mu\) be a strict associative chiral operation on \(A\),
and let \(u:\omega_X\to A\) be any differential-module map.
There are unique endomorphisms \(L,R\) with
\[
 \mu(u,1)=L\varepsilon_A,
 \qquad \mu(1,u)=-R\varepsilon_A\tau.
\]
They satisfy
\[
 \begin{gathered}
 L\mu=R\mu=0,\qquad L^2=LR=RL=R^2=0,\\
 \mu(L,1)=\mu(R,1)=\mu(1,L)=\mu(1,R)=0.
 \end{gathered}
\]
Normalization on either side forces \(A=0\).
\end{theorem}

\begin{proof}
The one-input case of Lemma~\ref{e1f:lem:insert} gives \(L,R\).
Insert \(u\) first in the associativity equation. The two
sides become
\[
 E_1\mu(L,1),\qquad E_1(L\mu)+E_2(L\mu).
\]
Independence of the two summands gives
\(L\mu=0=\mu(L,1)\). Inserting \(u\) last gives
\(R\mu=0=\mu(1,R)\); inserting it in the middle gives
\(\mu(R,1)=0=\mu(1,L)\). Precompose \(L\mu=R\mu=0\)
with \(u\) in each input. The one-input insertion isomorphism
then gives the four square-zero identities. A normalized action
has \(L=1_A\) or \(R=1_A\), so \(1_A=0\).
\end{proof}

The same argument controls maps to the chiral commutator.
Let \(f\) be a chiral Lie map from an algebra with vacuum to
\((A,\mu-\mu\tau)\), and form \(L,R\) using the image of
that vacuum. Preservation of its normalized action gives
\(f=(L+R)f\). But \((L+R)^2=0\), so \(f=0\).
Thus the universal associative recipient of a chiral algebra
with vacuum is zero in this ordinary pseudo-tensor category.

\begin{example}\label{e1f:ex:heisenberg}
At level \(\kappa\in k\), realize the Heisenberg vertex
algebra on \(k[x_1,x_2,\ldots]\). For \(h=x_1\), write
\(h(z)=\sum_nh_nz^{-n-1}\), where \(h_{-r}\) multiplies
by \(x_r\), \(h_r=\kappa r\partial_{x_r}\) for \(r>0\),
and \(h_0=0\). Translation is \(Tx_r=rx_{r+1}\).
The normal product \(a\cdot b=a_{(-1)}b\) satisfies
\[
 (h\cdot h)\cdot h-h\cdot(h\cdot h)=\kappa T^2h.
\]
Indeed, the constant coefficient of \(:h(z)^2:\) acting on
\(x_1\) is \(x_1^3+2h_{-3}h_1x_1=x_1^3+2\kappa x_3\),
whereas \(h_{-1}x_1^2=x_1^3\) and \(T^2h=2x_3\).
At zero level the normal product becomes polynomial
multiplication. At every level it is a regular coefficient
of a field, rather than a chiral map with localized input and
diagonal output. The vacuum obstruction applies at every level
to an associative operation of the latter type.
\end{example}

\subsection{Homotopies on the original coefficients}

For\phantomsection\label{e1f:def:termwise-homotopy-input} a complex \(A\), one can first form chiral Hom complexes
term by term. Their differential will be denoted \(\partial\).
Let \(\mu\) and \(u:\omega_X\to A\) be closed maps of
degree zero, and suppose a ternary operation of degree \(-1\)
satisfies
\[
 \partial m_3=\mu(\mu,1)-\mu(1,\mu).
\]
The insertion lemma remains an isomorphism of these termwise
complexes.

\begin{proposition}\label{e1f:prop:homotopy-vacuum}
Write \(\mu(u,1)=L\varepsilon_A\). There is a degree-minus-one
endomorphism \(K\) with \(\partial K=-L^2\). If
\(\partial a=L-1_A\), then
\[
 h=-K-a(L+1_A),\qquad \partial h=1_A.
\]
Thus a left vacuum unit up to homotopy makes \(A\) contractible.
If \(L=1_A\) and \(m_3(u,-,-)=0\), then \(A=0\).
\end{proposition}

\begin{proof}
Decompose
\(m_3(u,-,-)=E_1H_1+E_2H_2\) by the insertion lemma.
Comparing the two independent components gives
\[
 \partial H_1=\mu(L,1)-L\mu,
 \qquad \partial H_2=-L\mu.
\]
Define \(K\) by \(H_2(u,1)=K\varepsilon_A\). Since \(u\)
is closed, insertion commutes with the differential, and hence
\((\partial K)\varepsilon_A=-L^2\varepsilon_A\).
The one-input case of the lemma proves \(\partial K=-L^2\).
As \(L\) is closed,
\[
 \partial\bigl(-K-a(L+1_A)\bigr)
 =L^2-(L-1_A)(L+1_A)=1_A.
\]
If the specified value of \(m_3\) is zero, then \(H_2=K=0\).
Together with \(L=1_A\), this gives \(1_A=0\).
\end{proof}

The commutator has the corresponding homotopy obstruction.
Define \(R\) by \(\mu(1,u)=-R\varepsilon_A\tau\).
If \((L+R)\varepsilon_A\) is homotopic to \(\varepsilon_A\),
the insertion isomorphism on \(H^0\) gives
\([L]+[R]=[1_A]\). The proof of
Theorem~\ref{e1f:thm:vacuum}, applied in that \(H^0\)
pseudo-tensor category, gives
\([L]^2=[L][R]=[R][L]=[R]^2=0\). Thus \([1_A]=0\),
which is a contracting homotopy of \(A\).

A contractible complex need not be the zero complex; the
two-term identity complex is an example. This distinction
separates a strict vacuum obstruction from a homotopy vacuum
obstruction. Neither permits nonzero state cohomology.

\subsection{Derived operations on modules}

Derived Hom changes the operation complexes, so its degree
range must be checked separately.

\begin{proposition}[Derived operations on degree-zero coefficients]
\label{e1f:prop:derived-heart}
For right differential modules \(M_i,N\) in degree zero, put
\[
 \mathbb P_I^{\mathrm{ch}}(\{M_i\};N)
 =\operatorname{RHom}_{\mathcal D_{X^I}}
 \left(j_{I*}j_I^*(\boxtimes_iM_i),\Delta_{I*}N\right).
\]
Then \(H^q\mathbb P_I^{\mathrm{ch}}=0\) for \(q<0\), and
\(H^0\mathbb P_I^{\mathrm{ch}}=P_I^{\mathrm{ch}}\).
An \(A_\infty\) structure on a degree-zero module using these
operations has a strictly associative ordinary binary map.
Its two nested ternary products vanish. A local vacuum
normalized on either side, including normalization up to
homotopy, forces the module to be zero.
\end{proposition}

\begin{proof}
On a smooth curve the complement of the diagonals is locally
the inversion of one product of equations. Its open direct
image is exact. External products over \(k\) and closed
diagonal direct image are exact as well, so the source and
target are modules. An injective resolution of the target in
nonnegative degrees proves the two cohomology assertions.

The associativity homotopy makes the two binary composites
equal in \(H^0\), which is their ordinary chiral Hom space.
Theorem~\ref{e1f:thm:twostep} makes both zero. The same
\(H^0\) identification turns homotopy normalization into
ordinary normalization, and
Theorem~\ref{e1f:thm:vacuum} applies.
\end{proof}

\begin{remark}\label{e1f:rem:derived}
Positive Ext groups can still occur. In one normal coordinate
\(y\), put \(D=k\langle y,\partial_y\rangle\), with
\([\partial_y,y]=1\), and take the right modules
\(\omega=D/\partial_yD\), \(\delta=D/yD\).
The free resolution
\(0\to D\xrightarrow{\partial_y\cdot}D\to\omega\to0\)
gives the Hom complex
\(\delta\xrightarrow{\cdot\partial_y}\delta\).
Its differential is injective with one-dimensional cokernel.
Consequently \(\operatorname{Hom}_D(\omega,\delta)=0\) but
\(\operatorname{Ext}^1_D(\omega,\delta)=k\).
The derived insertion map therefore need not be an equivalence.

This class has degree \(1\), whereas an associativity
homotopy has degree \(-1\). More generally, the operation
spaces of a degree-zero diagonal insertion are discrete;
Theorem~\ref{oco:thm:rigidity} identifies every coherent
associative structure there with a strict one. Nontrivial
associativity homotopies require a different coefficient
complex, whose operations must be specified on that complex.
\end{remark}

\section{Binary coefficients and their associativity equations}
\label{e1c:sec:coefficients}

The vanishing of triple composites does not force the binary
operation to vanish. To describe the surviving operations, we
recover a chiral map from its normal derivative coefficients.
This also separates its regular field coefficient from its
singular coefficients.

Choose\phantomsection\label{e1c:def:left-coefficients} a coordinate \(x\) on an open subset of \(X\), and put
\(G=A\otimes\omega_X^{-1}\). This is the left-module form of
\(A\). Write \(T=\nabla_{\partial_x}\), so
\[
 T(fa)=fTa+f'a,
 \qquad (a\,dx)\partial_x=-(Ta)\,dx.
\]
After suppressing the output density, the transfer module on
the small diagonal has the polynomial description
\begin{equation}\label{e1c:eq:normal}
 Q_n(G)=G[\lambda_1,\ldots,\lambda_n]
     /(T+\lambda_1+\cdots+\lambda_n)G[\lambda_1,\ldots,\lambda_n].
\end{equation}
The relation is a relation of \(k\)-modules: \(T\) acts on
coefficients and is a connection, not an \(\mathcal O_X\)-linear
endomorphism. Eliminating \(\lambda_n\) gives
\(G[\lambda_1,\ldots,\lambda_{n-1}]\).
For\phantomsection\label{e1c:def:binary-transfer-coordinates} two inputs, use \(u=x_1-x_2\) and \(s=\lambda_1\).
Right differentiation in the two coordinates acts by \(s\)
and \(-T-s\); multiplication by \(u\) acts by \(\partial_s\).

\begin{proposition}\label{e1c:prop:binary}
A binary chiral map on \(A\) is equivalent, in this coordinate,
to \(k\)-bilinear maps on local sections
\(m:G\otimes_kG\to G\) and
\(c_s:G\otimes_kG\to G[s]\), compatible with restriction.
Write
\[
 c_s(a,b)=\sum_{n\geq0}\frac{s^n}{n!}c_n(a,b),
 \qquad
 B_s(a,b)=m(a,b)+\int_0^s c_t(a,b)\,dt.
\]
The polynomial \(c_s(a,b)\) is finite for each pair of inputs.
The required identities are
\begin{align}
 Tm(a,b)&=m(Ta,b)+m(a,Tb),\notag\\
 c_s(Ta,b)&=-s c_s(a,b),&
 c_s(a,Tb)&=(T+s)c_s(a,b),\label{e1c:eq:sesqui}\\
 B_s(a,fb)&=fB_s(a,b),&
 B_s(fa,b)&=\sum_{j\geq0}\frac{f^{(j)}}{j!}
                         \partial_s^jB_s(a,b).
 \label{e1c:eq:functions}
\end{align}
If \(F(a,b;g)\) denotes the value on
\(g(u)a\,dx_1\boxtimes b\,dx_2\), the reconstructed map is
\begin{align}
 F(a,b;u^q)&=\partial_s^{q+1}B_s(a,b)
                      &&(q\geq0),\label{e1c:eq:positive}\\
 F(a,b;u^{-r-1})&=\frac1{r!}\sum_{j=0}^r\binom rj
                         s^{r-j}B_s(T^ja,b)
                      &&(r\geq0).\label{e1c:eq:negative}
\end{align}
No finite-generation or flatness condition on \(A\) is imposed.
\end{proposition}

\begin{proof}
The Weyl relations of a differential-module map read
\[
 F(a,b;ug)=\partial_sF(a,b;g),
 \qquad
 F(a,b;\partial_{x_1}g)=-F(Ta,b;g)-sF(a,b;g).
\]
Start with \(B_s=F(a,b;u^{-1})\). Multiplication by \(u\)
recovers every nonnegative power and gives \(c_s=F(a,b;1)\).
Differentiation of \(u^{-r-1}\) then yields
\eqref{e1c:eq:negative} by induction. Differentiating the
regular test \(1\) gives \(c_s(Ta,b)=-sc_s(a,b)\).
Simultaneous translation gives
\(TB_s=B_s(Ta,b)+B_s(a,Tb)\); its constant term and its
derivative give the remaining identities in
\eqref{e1c:eq:sesqui}.

The output is rooted at \(x_2\). Thus \(f(x_2)\) multiplies
its coefficient, and \(f(x_1)\) acts by the Taylor operator
\(f(x_2+\partial_s)\). This is
\eqref{e1c:eq:functions}; its sum terminates on the polynomial
normal derivative coefficient.

Conversely, define \(F\) by the two displayed reconstruction
formulas. Let \(T_1\) act on its first input. For \(r\geq1\),
\[
 \partial_s(s+T_1)^rB_s
 =r(s+T_1)^{r-1}B_s+(s+T_1)^rc_s
 =r(s+T_1)^{r-1}B_s,
\]
since \((s+T_1)c_s=0\). This proves the multiplication
relation on negative powers. At \(u^{-1}\) it is
\(\partial_sB_s=c_s\), and on nonnegative powers it follows
from \eqref{e1c:eq:positive}. The differentiation recurrence
proves the second Weyl relation on negative powers; differentiating
\(c_s(Ta,b)=-sc_s(a,b)\) proves it on nonnegative powers.
Simultaneous translation supplies the other coordinate derivative.

The function identities are preserved by these recurrences.
For the first input, the \(f'\) in \(T(fa)\) cancels the
commutator of \(s\) with the Taylor operator of \(f\).
Thus the formulas respect the defining relations of the
localized right differential module. This argument uses those
relations directly and applies equally to modules with torsion.
\end{proof}

The corresponding field expansion is
\begin{equation}\label{e1c:eq:series}
 \begin{aligned}
 Y(a,u)b&=\sum_{n\geq0}c_n(a,b)u^{-n-1}
       +\sum_{r\geq0}\frac{m(T^ra,b)}{r!}u^r,\\
 F(a,b;g)&=\operatorname{Res}_{u=0}
                   e^{su}g(u)Y(a,u)b\,du.
 \end{aligned}
\end{equation}
For a Laurent test with finite lower bound, only finitely many
coefficients contribute to this residue. The formula constructs
a binary chiral operation; locality, associativity of fields
and a vacuum remain additional identities.
This is the local differential-module form of the chiral
endomorphism operad.\footnote{B.~Bakalov, A.~De Sole,
R.~Heluani and V.~G.~Kac, \emph{An operadic approach to vertex
algebra and Poisson vertex algebra cohomology}, Japan. J. Math.
\textbf{14} (2019), 249--342, \S6 and Appendix~A,
\href{https://arxiv.org/abs/1806.08754}{arXiv:1806.08754}.}

To compare a second coordinate \(w=h(x)\), apply the recovered
operation to
\[
 \frac{h'(x_1)h'(x_2)}{h(x_1)-h(x_2)}
                    a\,dx_1\boxtimes b\,dx_2.
\]
Express the answer with \(dw=h'(x)dx\) and
\(\partial_w=(h')^{-1}\partial_x\) in the transfer module.
It must equal the polynomial assigned in the \(w\)-coordinate.
Only finitely many normal derivatives occur on each input.
Substitution and transfer compose, so these overlap conditions
satisfy the cocycle identity and glue the local maps.

\subsection{Computing the two ternary channels}

The translation of an inner output acts before that output
enters the outer map. Keeping it in that order, define
\begin{equation}\label{e1c:eq:H}
 H_{s,t}(a,b)=(T+s+t)m(a,b)
 +\sum_{n\geq0}
   \frac{s^{n+2}-(-T-t)^{n+2}}{(n+2)!}c_n(a,b).
\end{equation}
Here every \(T\) acts on the displayed inner coefficient.
Equivalently, this polynomial evaluates the formal integral
\(\int_0^{T+s+t}B_{s-\tau}(a,b)\,d\tau\).

\begin{theorem}\label{e1c:thm:criterion}
The binary map of Proposition~\ref{e1c:prop:binary} is strictly
associative exactly when, for every \(a,b,c\),
\begin{equation}\label{e1c:eq:criterion}
 \begin{aligned}
 &B_{\lambda_1+\lambda_2}
       \bigl(H_{\lambda_1,\lambda_2}(a,b),c\bigr)\\
 &\hspace{18mm}
   +B_{\lambda_1}\bigl(a,H_{\lambda_2,\lambda_3}(b,c)\bigr)=0
                      \quad\text{in }Q_3(G).
 \end{aligned}
\end{equation}
Each summand then vanishes separately. In this formula, expand
the inner \(H\), apply the outer \(B\) to its coefficients,
and only then set
\(\lambda_3=-T_{\mathrm{out}}-\lambda_1-\lambda_2\).
\end{theorem}

\begin{proof}
Use\phantomsection\label{e1c:def:ternary-scalar-test} \(u=x_1-x_2\), \(v=x_2-x_3\), and the scalar input
\(\Omega=1/(uv(u+v))\). For the first collision its remaining
factor is \(1/(v(v+u))\). An inner polynomial \(P(s)\)
of normal derivatives acts on that factor by
\[
 \left.P(s+\partial_u)\frac1{v(v+u)}\right|_{u=0}
 =\sum_{j\geq0}\frac{(-1)^jP^{(j)}(s)}{v^{j+2}}
 =\left(\int_0^{-\partial_v}P(s-\tau)\,d\tau\right)\frac1v.
\]
The transfer relation moves \(-\partial_v\) to
\(T+\lambda_1+\lambda_2\) on the first input of the outer
operation. Integrating the monomials of \(B_s\) gives
\[
 (\mu\circ_1\mu)(a,b,c;\Omega)
 =B_{\lambda_1+\lambda_2}
             \bigl(H_{\lambda_1,\lambda_2}(a,b),c\bigr).
\]

For the other collision, keep \(t=x_1-x_3\) and
\(v=x_2-x_3\). The remaining factor is \(1/(t(t-v))\),
and the corresponding identity is
\[
 \left.P(s+\partial_v)\frac1{t(t-v)}\right|_{v=0}
 =\sum_{j\geq0}\frac{P^{(j)}(s)}{t^{j+2}}
 =-\left(\int_0^{\partial_t}P(s-\tau)\,d\tau\right)\frac1t.
\]
After transfer this gives
\[
 (\mu\circ_2\mu)(a,b,c;\Omega)
 =-B_{\lambda_1}\bigl(a,H_{\lambda_2,\lambda_3}(b,c)\bigr).
\]
The minus sign accounts for the sum in
\eqref{e1c:eq:criterion}.

This one scalar test generates all tests under multiplication
and differentiation. Indeed, for \(a,b>0\),
\[
 \frac1{u^av^b(u+v)^c}
 =\frac1{u^{a-1}v^b(u+v)^{c+1}}
  +\frac1{u^av^{b-1}(u+v)^{c+1}}.
\]
Iterating reduces to denominators using two independent linear
forms. Their simple-pole products are polynomial multiples of
\(\Omega\); derivatives in the dual coordinate directions
produce higher poles. Multiplication then includes single-pole
and regular tests. For arbitrary input sections, commute each
scalar derivative past those sections and subtract the terms
where it differentiates an input. Induction proves the same
generation assertion without freeness assumptions on \(G\).
Thus equality on \(\Omega\) is equality of the ternary maps.
Theorem~\ref{e1f:thm:twostep} makes their values vanish
separately, and the converse is immediate.
\end{proof}

These transfer formulas agree with those of Dudin and
Kolesnikov after \(z_i=-x_i\) and \(\partial=T\).\footnote{
I.~V.~Dudin and P.~S.~Kolesnikov, \emph{Chiral algebras with
abelian conformal part},
\href{https://arxiv.org/abs/2606.10367}{arXiv:2606.10367v1},
equations~(13), (16). Their Theorem~1 classifies the regular
case for binary quadratic operads.}
For the scalar residue on horizontal inputs,
\(m(1,1)=1\) and \(c_s=0\). The two composites on
\(\Omega\) are \(\lambda_1+\lambda_2\) and \(\lambda_1\).
Their difference \(\lambda_2\) retains a normal derivative;
projecting immediately to a scalar residue would lose it.

\subsection{Regular and singular nilpotent operations}

Call\phantomsection\label{e1c:def:regular-binary-input} a binary chiral map \emph{regular} if it vanishes on the
unlocalized input \(A\boxtimes A\). This is the condition
\(c_s=0\). Its field series is regular, although its values
on polar inputs can be nonzero.

\begin{corollary}\label{e1c:cor:regular}
Regular strict associative chiral operations are equivalent
to \(\mathcal O_X\)-bilinear products
\(m:G\otimes_{\mathcal O_X}G\to G\)
for which the connection is a derivation and
\[
 m(m(a,b),c)=0,
 \qquad m(a,m(b,c))=0.
\]
Neither commutativity nor anticommutativity is required.
\end{corollary}

\begin{proof}
For \(c_s=0\), the function relations make \(m\)
\(\mathcal O_X\)-bilinear. The first summand of the criterion is
\[
 m(Tm(a,b),c)+(\lambda_1+\lambda_2)m(m(a,b),c).
\]
Eliminating the final translation makes the other summand
\(-m(Ta,m(b,c))-\lambda_1m(a,m(b,c))\).
Its \(\lambda_2\) and \(\lambda_1\) coefficients give the
two identities successively. Conversely, applying those
identities also to differentiated inputs kills the constant
terms and proves the criterion.
\end{proof}

For\phantomsection\label{e1c:def:three-generator-nilpotent-example} example, take horizontal generators \(p,q,z\) and the
single nonzero product \(m(p,q)=z\). Tensoring this product
with the scalar principal-parts operation gives a nonzero
chiral multiplication: a simple pole on \(p,q\) maps to
\(z\) on the diagonal. Every further product vanishes because
\(z\) annihilates both inputs.

Transposition has a sign different from the ordinary opposite
product. In these coordinates,
\[
 B_s^{\mathrm{op}}(a,b)=-B_{-T-s}(b,a).
\]
The minus sign reverses the normal coordinate; for homogeneous
super inputs the right side also has the Koszul factor
\((-1)^{|a||b|}\). In the regular ungraded case the opposite
normal product is \(-m(b,a)\). Consequently the chiral
commutator has normal product \(m(a,b)+m(b,a)\), and its
kernel among these associative operations is exactly the
antisymmetric normal products.

Every symmetric differential product \(n\) with zero nested
triples is the chiral commutator of \(m=n/2\). Its residue
operation is Lie: reversing the two positions supplies the
skew sign, and each ternary term is zero. Conversely, operadic
symmetry \(\mu=\mu^{\mathrm{op}}\) in the regular sector
requires an antisymmetric normal product. It is therefore
different from a commutative differential algebra, whose full
field product becomes a chiral Lie operation after density
conversion.

Singular\phantomsection\label{e1c:def:singular-two-step-sum} examples have the same annihilation mechanism.
Given differential modules \(P,Z\) and a chiral map
\(q:j_*j^*(P\boxtimes P)\to\Delta_*Z\), set \(A=P\oplus Z\).
Use \(q\) on the two \(P\) inputs and zero whenever an input
lies in \(Z\). Both nested triples vanish. If \(q\) is
nonzero on regular inputs, this gives a singular strictly
associative operation.

For\phantomsection\label{e1c:def:singular-coefficient-ideal} a general strict associative operation, let \(J\) be the
smallest differential submodule containing all \(c_n(a,b)\)
and stable under \(m\) and every \(c_n\) in either input.

\begin{corollary}\label{e1c:cor:quotient}
The quotient \(G/J\) has an induced differential product
with zero nested triples. Every normal-product expression
on at least three inputs belongs to \(J\).
\end{corollary}

\begin{proof}
Stability makes all coefficient operations descend. On the
quotient the \(c_n\)'s vanish, so the regular classification
applies. Under coordinate change, the new singular coefficients
are combinations of the old ones. The change in the normal
coefficient is also a combination of singular coefficients,
by the finite Taylor overlap formula. Applying the inverse
coordinate change shows equality of the resulting submodules.
Thus \(J\) is intrinsic and glues on \(X\).
\end{proof}

The normal coefficient cannot have an ordinary two-sided unit
on a nonzero module. If \(e\) were one, the derivation identity
would give \(Te=T m(e,e)=2Te\), hence \(Te=0\).
Sesquilinearity would then give
\(s c_s(e,a)=0\) and \((T+s)c_s(a,e)=0\).
Both polynomial operators are injective, the second by its
leading coefficient. Thus both singular actions vanish, and
\(e\) defines the normalized residue action excluded by
Theorem~\ref{e1f:thm:vacuum}.

\subsection{Unbounded poles prevent an ordinary envelope}

Removing the vacuum permits the preceding examples, but does
not give an associative envelope functor on all differential
modules. A commutator fixes the skew part of the operation
and leaves its symmetric pole order uncontrolled.

\begin{proposition}\label{e1f:prop:noenvelope}
For \(X=\mathbb A^1\), the commutator functor from associative
to Lie algebras in \(P_X^{\mathrm{ch}}\), on all right
\(\mathcal D_X\)-modules, has no left adjoint.
\end{proposition}

\begin{proof}
Let \(M=x\mathcal D_X\oplus y\mathcal D_X\) have zero
bracket; \(x,y\) here name free module generators.
Use\phantomsection\label{e1f:def:envelope-density} \(t=z_1-z_2\) and \(\eta=dz_1\wedge dz_2\).
For each \(r>0\), prescribe on the generators the map
\[
 \beta_r(fx\boxtimes y)
       =\operatorname{pp}_{t=0}(ft^{-r}\eta)
                   \quad\text{in }\Delta_*\omega_X,
\]
and extend by differential-module induction. On
\(A_r=M\oplus\omega_Xc\), use \(\beta_r\) for the ordered
pair \(x,y\), its transpose for \(y,x\), and zero on all
other generator pairs and on every input from \(\omega_Xc\).
Every inner output annihilates further inputs, so the resulting
operation is associative. It is symmetric, hence has zero
chiral commutator; the inclusion \(M\to A_r\) is a Lie map.

If a universal envelope \(U\) existed, the section
\(w=\mu_U(x\boxtimes y)\) of \(\Delta_*U\) would be
killed by some \(t^n\). Its image under the algebra map
\(U\to A_r\) extending the inclusion of \(M\) would be
the principal part of \(t^{-r}\eta\). For \(r>n\), that
image is not killed by \(t^n\), a contradiction.
\end{proof}

The\phantomsection\label{e1f:def:unbounded-symmetric-central-product} obstruction persists with a fixed nonzero central bracket.
Transposition sends \(t\) to \(-t\) and \(\eta\) to
\(-\eta\), so it acts on \(\beta_r\) by \((-1)^{r+1}\).
On \(a\mathcal D_X\oplus\omega_Xc\), make \(c\)
annihilate both inputs and prescribe
\[
 \mu_{\kappa,r}(a,a)=\left(\frac\kappa2\beta_2+\beta_r\right)c
                       \qquad(r\geq3\text{ odd}).
\]
All nested triples vanish, and all these maps have commutator
\(\kappa\beta_2c\). Their symmetric pole orders are unbounded,
so the same argument rules out an envelope of this fixed
nonunital central Lie algebra. It is not the Heisenberg vertex
algebra with vacuum.

The geometric envelope of Proposition~\ref{cqe:prop:envelope}
lives instead in the symmetric monoidal collision category.
It retains configuration coefficients and an empty unit.
Its existence therefore asserts no left adjoint on the
ordinary curve modules used here.

\section{Associative objects in vertex algebras}
\label{e1f:sec:va}

One can ask for an associative multiplication while keeping
the full vertex structure. Its source is then the ordinary
tensor product of vertex algebras, rather than the localized
input of a chiral operation. This change avoids the preceding
definition but imposes a different restriction.

\phantomsection\label{e1f:def:regular-vertex-algebra}A vertex algebra has a state space \(V\), a vacuum
\(\mathbf1\), a translation \(T\), and fields with
\(Y(a,z)b\in V((z))\). They obey
\[
 \begin{gathered}
 Y(\mathbf1,z)=1_V,\qquad Y(a,z)\mathbf1=e^{zT}a,
 \qquad T\mathbf1=0,\\
 [T,Y(a,z)]=\partial_zY(a,z)=Y(Ta,z).
 \end{gathered}
\]
For each pair of states, locality says that a power of
\(z-w\) kills the commutator of their fields. Call the vertex
algebra \emph{regular} if every \(Y(a,z)b\) belongs to
\(V[[z]]\).

\begin{theorem}\label{e1f:thm:va}
Unital associative objects for the ordinary tensor product
of vertex algebras are equivalent to unital commutative
algebras with derivation. Their fields are
\[
 Y(a,z)b=(e^{zT}a)b,
\]
and in particular have no singular terms.
\end{theorem}

\begin{proof}
Let \(m:V\otimes V\to V\) be such a multiplication.
Its unit is the vertex vacuum, since the monoidal unit is
the scalar vertex algebra. Applying preservation of fields
to \(a\otimes\mathbf1\) and \(\mathbf1\otimes b\) gives
\[
 \begin{aligned}
 Y(a,z)b
 &=Y(m(a,\mathbf1),z)m(\mathbf1,b)\\
 &=m\bigl(Y(a,z)\mathbf1,Y(\mathbf1,z)b\bigr)
  =m(e^{zT}a,b).
 \end{aligned}
\]
Using \(\mathbf1\otimes a\) and \(b\otimes\mathbf1\)
instead gives \(Y(a,z)b=m(b,e^{zT}a)\). The constant
coefficient therefore makes \(m\) commutative.
Preservation of translation makes \(T\) a derivation.

Conversely, a unital commutative algebra with derivation has
the displayed regular fields. Its translation exponential
preserves multiplication, so \(m\) preserves these fields
and defines the required associative object. A map in either
description preserves exactly multiplication, unit and
translation, which proves the equivalence of categories.
\end{proof}

The unit has forced every field value to be a translate of
the same-position multiplication. Singular field products
therefore require a source allowing distinct positions.
The strict rectangular comparison makes precise what an
ordinary associative object recovers under an independently
specified state functor. For its singleton normalization,
write \(H_V=\omega_X\otimes_kV\) for the right differential
module of a translation-equivariant state space \(V\).

\input{manuscript/constructions/strict-rectangular-comparison.tex}

\section{Ordered products with matrix coefficients}
\label{e1f:sec:ordered}

Associativity\phantomsection\label{e1f:def:matrix-field-inputs} of ordered fields does not require the binary
product itself to be a vertex-algebra morphism. Let \(V\)
be a nonzero vertex algebra and \(B\) a unital associative
algebra. On \(V\otimes B\), define
\begin{equation}\label{e1f:eq:VB}
 Y(a\otimes b,z)(a'\otimes b')
                 =Y_V(a,z)a'\otimes bb'.
\end{equation}
The vacuum is \(\mathbf1_V\otimes1_B\) and translation is
\(T_V\otimes\mathrm{id}_B\). Every iterated field product is the
corresponding product in \(V\), multiplied by the coefficient
word \(b_1\cdots b_n\) in its original order. Associativity
of \(B\) identifies the expansions from different consecutive
bracketings. Reordering inputs requires a further operation.

\begin{proposition}\label{e1f:prop:nonlocal}
The fields in \eqref{e1f:eq:VB} are mutually local exactly
when \(B\) is commutative.
\end{proposition}

\begin{proof}
Write\phantomsection\label{e1f:def:left-multiplication-operator} \(L_P\) for left multiplication by \(P\in B\).
The field of \(a\otimes P\) is \(Y_V(a,z)\otimes L_P\).
If the \(L_P\)'s commute, locality follows from that of \(V\).
For the states \(\mathbf1_V\otimes P\) and
\(\mathbf1_V\otimes Q\), their field commutator is the
constant operator \(1_V\otimes L_{[P,Q]}\). No nonzero
power of \(z-w\) annihilates it. Locality therefore implies
\(L_{[P,Q]}=0\), and evaluation at \(1_B\) gives
\([P,Q]=0\).
\end{proof}

Take \(P=E_{12}\), \(Q=E_{21}\) in \(\operatorname{Mat}_2(k)\).
Their commutator is \(\operatorname{diag}(1,-1)\), so this
failure persists at separated positions. It cannot be encoded
by a map whose output is supported only on the collision
diagonal.

Nonconsecutive collisions retain another part of the order.
For \(V=H_\kappa\), with \(\kappa\ne0\), consider the
ordered fields of \(h\otimes P\), \(\mathbf1\otimes Q\),
and \(h\otimes P\). Contracting the outer Heisenberg fields
gives the coefficient
\[
 \frac{\kappa PQP}{(z_1-z_3)^2}
       =\frac{\kappa P}{(z_1-z_3)^2},
 \qquad P^2=0.
\]
Replacing the two outer coefficients by their product would
erase this pole. The middle matrix must remain in the
coefficient of the partial collision.

Call a collection a \emph{simultaneous ordered product system}
when each finite ordered input has a single multilinear product
with a finite common denominator on its collision diagonals,
and every consecutive substitution tree is an expansion of
that product. Require the vacuum and translation identities
stated above. The systems \eqref{e1f:eq:VB} have these
properties. In their regular sector, all positions may be set
equal without taking a principal part.

\begin{proposition}\label{e1f:prop:regular}
Regular simultaneous ordered product systems on \(\mathbb A^1\)
are equivalent to unital associative algebras with derivation,
through \(Y(a,z)b=(e^{zT}a)b\). The system is local precisely
when the algebra is commutative.
\end{proposition}

\begin{proof}
Define \(ab=Y(a,0)b\). Regularity lets one cancel the
regularizing factor in weak associativity and evaluate at
zero; the resulting identity is \((ab)c=a(bc)\).
The two vacuum identities give the unit. Translation
covariance gives the derivation identity and determines every
Taylor coefficient from \(\partial_zY(a,z)=Y(Ta,z)\).
Thus \(Y(a,z)b=(e^{zT}a)b\).

Conversely, ordered products of translated inputs give the
simultaneous products of an associative differential algebra.
The derivation law and associativity give every consecutive
substitution identity. If the fields are local, cancellation
of the locality factor in their regular commutator followed
by evaluation on the vacuum at zero gives \(ab=ba\).
Commutativity conversely makes all translated products local.
\end{proof}

An exchange would relate two input orders by a spectral
operator on their states. Regularity and associativity do not
provide such an operator with finite state coefficients.

\begin{proposition}\label{e1f:prop:noexchange}
Let \(B\) be free associative on \(a_i,b_i\), \(i\geq0\),
with \(Ta_i=a_{i+1}\), \(Tb_i=b_{i+1}\). There is no
expression
\[
 (e^{uT}a_0)b_0
       =\sum_{\nu=1}^N f_\nu(u)c_\nu(e^{uT}d_\nu)
\]
with \(f_\nu(u)\in k((u))\) and fixed \(c_\nu,d_\nu\in B\).
Thus its regular ordered system has no exchange of this finite
Laurent-series type.
\end{proposition}

\begin{proof}
Project onto the span of \(a_ib_j\), identified with
\(k[p,q]\) by \(a_ib_j\mapsto p^iq^j\). The left side
becomes \(e^{up}\). On the right, the translated suffix can
be empty, the final factor, or the entire word. Each term
therefore has the form
\(f(u)P(p,q)e^{uL}\), with
\(L\in\{0,q,p+q\}\). Let \(d\) bound the degrees of
the finitely many polynomials \(P\), and put
\(D=\partial_p+2\partial_q\).
The commuting operator
\[
 \prod_{s\in\{0,2,3\}}(D-su)^{d+1}
\]
annihilates every term on the right. On the left it gives
\((2u^3)^{d+1}e^{up}\), which is nonzero. This contradicts
the proposed expression.
\end{proof}

The conclusion concerns its specified finite coefficient
domain. A completion allowing other state sums would require
its own exchange and substitution maps. Even with no poles
in the fields, the coefficient topology is part of that
construction.

\section{Words, cuts and associative chiral algebras}
\label{e1f:sec:tensor}

The coefficient \(PQP\) in the outer Heisenberg collision
records an order even when the first and last positions agree.
The set of occupied positions alone forgets that order.
We therefore begin with words in positions and identify only
adjacent repetitions.

Let \(\mathsf{Surj}_{\mathrm{mon}}\) have the nonempty ordered
sets \([n]=\{1<\cdots<n\}\) and monotone surjections.
A surjection gives a diagonal map repeating each entry over
its interval fibre.

\begin{proposition}\label{e1f:prop:words}
For a set \(S\), the colimit of \(S^{[n]}\) under these
diagonal maps is the set of nonempty words modulo insertion
and deletion of adjacent equal entries. Every class has a
unique representative with no adjacent repetition.
For distinct \(a,b\), the words \(ab,aba,abab,\ldots\)
are distinct although their supports agree.
\end{proposition}

\begin{proof}
A monotone surjection repeats each entry a positive number
of consecutive times. Replacing every maximal run of equal
entries by one entry is invariant under these moves.
Every word is related to its reduction, and a reduced word
cannot change under a sequence of these moves followed by
reduction. This proves existence and uniqueness.
\end{proof}

The word \(aba\) is precisely an outer collision with a
middle position retained. On a curve, its coefficients belong
to the category \(\mathcal W_X\) of complexes on \(X^n\)
with coherent exceptional-pullback identifications under the
monotone diagonal maps. Adjoin an isolated empty coefficient
in \(D(k)\). For two systems, define the cut convolution by
\begin{equation}\label{e1f:eq:cut}
 (F\star G)_I
   =\bigoplus_{I=A|B}j_{A|B*}j_{A|B}^*(F_A\boxtimes G_B).
\end{equation}
The sum is over an initial interval \(A\) and its final
complement \(B\), including empty intervals. The open set
of \(j_{A|B}\) separates every position in \(A\) from every
position in \(B\).

Both bracketings of a triple convolution localize at the
same three consecutive blocks. Open base change identifies
them and gives the higher associativity identifications.
The unit \(K=k_\varnothing\) is supported on the empty word.
This tensor is not symmetric: point probes at distinct
\(a,b\) give binary products supported at \((a,b)\) and
\((b,a)\), respectively.

There is a second configuration category. Take a chart
\(X^I\) for each total order on \(I\), and glue charts
where the orders do not conflict
(Chapter~\ref{ogd:foundation}). Over a configuration whose
equality blocks are \(B\), the resulting ordered collision
scheme has \(\prod_B|B|!\) points. Its convolution sums over
ordered decompositions into arbitrary subsets, and is
symmetric monoidal. It has the same empty unit, but different
nonconsecutive collision maps from \(\mathcal W_X\).
The separate-variable tensor of Chapter~\ref{cl:chapter}
is a further coefficient category, with its own scalar unit.

\subsection{The finite bar and cobar constructions}

Call a system \emph{positive} if its empty coefficient is
zero. A product of \(r\) positive systems vanishes on a
configuration with fewer than \(r\) labels: every factor
requires a nonempty block. This holds for each convolution
just described and bounds the lengths of its normalized bar
and cobar words.

For\phantomsection\label{e1f:def:bar-augmentation-suspension} an augmented algebra \(A\to K\) with \(A_\varnothing=k\),
write \(\overline A\) for its augmentation ideal and
\(s\overline A=\overline A[1]\). The bar has underlying
coalgebra
\[
 BA=K\oplus\bigoplus_{r\geq1}(s\overline A)^{\star r},
\]
with deconcatenation coproduct. On a strict dg model use
\[
 b_1(sa)=-s(da),\qquad
 b_2(sa\otimes sb)=(-1)^{|a|}s(ab),
\]
and extend as a coderivation with tensor-composition signs.
Its square expresses the differential and associativity
identities of \(A\).

For a coaugmented coalgebra \(C=K\oplus\overline C\),
cobar is the tensor algebra on \(t\overline C=\overline C[-1]\).
Writing \(\overline\Delta c=\sum c'\otimes c''\), its
differential is the derivation determined by
\[
 d(tc)=-t(dc)+\sum(-1)^{|c'|+1}(tc')(tc'').
\]
For coherent algebras and coalgebras, the normalized bar and
cobar retain the higher operations. Every term remains finite
at each configuration.

Bar and cobar are inverse equivalences when the empty
coefficient is \(k\)
(Theorem~\ref{ogd:thm:duality} and
Section~\ref{ogd:sec:cut-support}). Their comparison can be
seen by flattening an iterated bar--cobar word. Each gap
between its original factors is either inside a block or
between blocks. On the associated graded, changing that
choice is an identity map between two consecutive degrees.
The complex of all gap choices is a tensor product of these
two-term identity complexes. It contracts unless there is
just one original factor. The empty term and that factor
survive and give the comparison map. The configuration bound
makes the filtration finite, even for unbounded coefficient
complexes. No completion or extra global elementwise
conilpotence assumption is needed here.

This is associative Koszul self-duality in the collision
category: \(\mathrm{Ass}^!=\mathrm{Ass}\) relates an algebra
to a coalgebra. Turning that coalgebra into a second algebra
requires an additional duality, such as the positive integer
weight dual of Proposition~\ref{cqe:prop:dual}.

\subsection{Factorization and the support of the cobar}

A coalgebra is \emph{factorizing} for the chosen convolution
if each specified coproduct component becomes an equivalence
to the external product on its separating open set. For cut
convolution this is required at consecutive cuts. On the
ordered collision schemes it is required for every ordered
decomposition into subsets.

\begin{definition}\label{e1f:def:e1}
For either specified convolution, an \emph{associative chiral
algebra} is an augmented associative algebra \(A\) whose
augmentation ideal is positive and whose bar \(BA\) is
factorizing.

A \emph{vacuum marking} consists of a left differential
module \(G\), a horizontal map \(\mathcal O_X\to G\), and
an identification
\(A_{\{1\}}\cong\omega_X\otimes_{\mathcal O_X}G\).
A \emph{vacuum framing} also specifies a comparison with
ordered simultaneous products on \(G\), preserving
substitutions, translation and that vacuum.
\end{definition}

On\phantomsection\label{e1f:def:cut-collision-support} words, the factorization condition has a direct geometric
meaning. Let \(U_p\subset X^n\) be the open set separating
the first \(p\) labels from the remaining labels, and put
\[
 Z_1=X,\qquad
 Z_n=X^n\setminus\bigcup_{p=1}^{n-1}U_p
     =\bigcap_{p=1}^{n-1}\ \bigcup_{i\leq p<j}\Delta_{ij}.
\]
Thus every consecutive cut at a point of \(Z_n\) is crossed
by an equality. In particular, \(Z_3=\Delta_{13}\): its
middle position may remain distinct.

For a strong cut coalgebra \(C\), iterated coproducts identify
its cobar coefficient with the augmented localization cube
of the opens \(U_p\). The nonempty vertices give the \v{C}ech
complex of their union; adjoining the empty vertex takes
the fibre of restriction to that union. With the cobar
desuspension, the result is
\[
 \overline{\Omega_{\mathrm{cut}}C}_n
                 \simeq R\Gamma_{Z_n}(C_n)[-1]
\]
(Theorem~\ref{ogd:thm:cut-local-cohomology}). Its
multiplication uses the inverse factorization maps on the
cut opens, in addition to the supported coefficient object.
Conversely, if \(\overline A_n\) is supported on \(Z_n\),
every surviving bar word on a cut open has a separator at
that cut. Deconcatenation there is an equivalence.
Thus \(BA\) is factorizing exactly when \(A\)'s positive
coefficients have these supports
(Theorem~\ref{ogd:thm:Z-restriction}).

\begin{remark}\label{e1f:rem:two-carriers}
Factorization on the ordered collision schemes has a
different support consequence. Let \(C=BA\), let \(U\)
be the two-point distinct locus, and put
\(V=(C_{\{1\}}\boxtimes C_{\{1\}})|_U\).
If both ordered coproduct components are equivalences,
one identifies their joint map with
\[
 (1,R):V\longrightarrow V\oplus V,
\]
where \(R\) compares the two orders. The quotient map
\((x,y)\mapsto y-Rx\) identifies its cofiber with \(V\).
Consequently
\[
 \overline A_{\{1,2\}}|_U
       \simeq V[-2]
       \simeq(\overline A_{\{1\}}
                         \boxtimes\overline A_{\{1\}})|_U.
\]
The two binary product maps are the two summand inclusions
in that cofiber, hence \(-R\) and \(1\). They are both
equivalences. If \(V\ne0\), the algebra is not supported
on the collision locus; \eqref{ogd:eq:U} gives an example.

By contrast, the bar of an algebra supported on the small
diagonal branches has, on \(U\), two singleton words.
Its coproduct components are the projections
\(V\oplus V\to V\), neither an equivalence when \(V\ne0\).
Their \emph{joint} map is the relevant equivalence.
Theorem~\ref{ogd:thm:joint-recognition} characterizes
small-diagonal support by this joint word factorization.
It cannot be replaced by the componentwise condition in
Definition~\ref{e1f:def:e1}.
\end{remark}

The two units now have their specified locations. The
associative unit comes from \(K\) on the empty configuration.
The local vacuum is a singleton state and may occur as the
output of a singular product. For the operation-level
formulation with spectators
(Definition~\ref{sop:def:algebra}), its normalized unit is
a morphism from the unit chiral algebra. Applying the residue
coalgebra and cobar constructions gives the vacuum marking
(Proposition~\ref{sop:prop:vacuum-morphism}).

Separating\phantomsection\label{e1f:def:degree-zero-diagonal-word} the units alone does not alter the obstruction
on a degree-zero module. Its diagonal word insertion \(iM\)
has \((iM)_n=\Delta_{n*}M\) and empty coefficient zero.
Theorem~\ref{oco:thm:operations} identifies the operations
from its cut powers to itself with the derived chiral
operation complexes. Their negative cohomology vanishes,
so every associative structure on \(K\oplus iM\) has the
strict two-step nilpotent product of
Theorem~\ref{oco:thm:rigidity}. A normalized local vacuum
still forces \(M=0\).
The field construction instead retains the complexes on
\(Z_n\). The scalar vacuum already shows why they are larger.

\section{Collision coefficients of the scalar vacuum}

Take\phantomsection\label{e1f:def:scalar-vacuum-field} the scalar field algebra \(k\), with translation zero
and \(Y(1,z)1=1\). Requiring its products to be regular
gives the coordinate lattice \(\mathcal K_I=\mathcal O_{X^I}\).
Its normalized coefficient coalgebra is
\begin{equation}\label{e1f:eq:scalar-coalgebra}
 C_I=(\omega_{X^I}\otimes_{\mathcal O_{X^I}}\mathcal K_I)[|I|]
                  =\omega_{X^I}[|I|],
 \qquad C_\varnothing=k.
\end{equation}
The shift makes exceptional collision pullback identify the
coefficient on \(I\) with that on the collapsed set
(Theorem~\ref{cl:thm:exceptional}). At one position it is
\(\omega_X[1]\); its cobar desuspension is \(\omega_X\).
The higher coefficients are determined by the cuts.

For\phantomsection\label{e1f:def:scalar-cut-ideal} \(X=\mathbb A^1\), let \(R=k[z_1,\ldots,z_n]\), and
put
\[
 F_a=\prod_{i\leq a<j}(z_i-z_j),
 \qquad \mathfrak a_n=(F_1,\ldots,F_{n-1}).
\]
The cut after \(a\) localizes at \(F_a\). A set \(S\)
of cuts gives \(|S|+1\) cobar blocks. Their density shifts
add to \([n]\), and each block is desuspended once.
After trivializing the density, the corresponding term is
\begin{equation}\label{e1f:eq:scalar-cobar-term}
 R\left[(\prod_{a\in S}F_a)^{-1}\right][n-|S|-1],
\end{equation}
in\phantomsection\label{e1f:def:scalar-cobar-coefficient} degree \(|S|-(n-1)\). The alternating localization maps
give the augmented \v{C}ech complex. If \(\mathcal L_n\)
denotes this reduced cobar coefficient, then
\[
 \mathcal L_n\simeq R\Gamma_{\mathfrak a_n}(R)[n-1].
\]
There are no cuts at one point, so \(\mathcal L_1=R\).
At two points the support is the diagonal.

At three points, write \(h=z_1-z_3\), \(s=z_1-z_2\),
and \(v=z_2-z_3\). Then \(F_1=sh\), \(F_2=hv\), and
\(F_1+F_2=h^2\). It follows that
\(\sqrt{\mathfrak a_3}=(h)\), and hence
\[
 H^{-1}(\mathcal L_3)=R[h^{-1}]/R.
\]
This coefficient lies on the outer diagonal with the middle
position free. At four points the support is
\[
 \{z_1=z_4\}\ \cup\ \{z_1=z_3,\ z_2=z_4\}.
\]
The second component can have two distinct collision positions.

Call a set of disjoint label pairs an \emph{irredundant
spanning matching} if its pairs cross every consecutive cut
and no pair can be removed while retaining that property.

\begin{theorem}\label{e1f:thm:support}
For \(n\geq2\), the irreducible components of
\(V(\mathfrak a_n)\) are
\[
 \{z_a=z_b\text{ for every }\{a,b\}\in M\},
\]
indexed by irredundant spanning matchings \(M\).
There are \(\binom{n-2}{2r-2}\) components of codimension
\(r\), for \(1\leq r\leq\lfloor n/2\rfloor\).
For \(n\geq3\) their total number is \(2^{n-3}\).
At the generic point of such a component, with independent
normal coordinates \(u_1,\ldots,u_r\), the cohomology of
\(\mathcal L_n\) is concentrated in degree \(r-n+1\)
and contains \(1/(u_1\cdots u_r)\), whose iterated normal
residue is \(1\).
\end{theorem}

\begin{proof}
The equations \(F_a=0\) say that each cut is crossed by
an equality. The equality blocks of a point cover cuts:
a block covers exactly those between its least and greatest
labels. Replacing it by the equality of those two labels
preserves this coverage and removes every other equation
of the block. More generally, an equality cluster formed
by pairs sharing labels can be replaced by its outer pair.
Thus a minimal set of coincidence equations consists of
disjoint pairs, each needed for at least one cut. These are
exactly the irreducible components of the union of equality
subspaces.

Write its pairs as \(\{a_i,b_i\}\), with \(a_i<b_i\),
ordered by their left endpoints. No interval
\([a_i,b_i)\) is contained in another, so the right
endpoints increase as well. Coverage of the first and last
cuts gives \(a_1=1\), \(b_r=n\). Consecutive intervals
overlap: \(a_{i+1}<b_i\), since the endpoints are distinct
and otherwise a cut is uncovered. But
\(b_i<a_{i+2}\), since overlap of intervals \(i,i+2\)
would make interval \(i+1\) redundant. Therefore their
endpoints have the order
\[
 1=a_1<a_2<b_1<a_3<b_2<\cdots
                  <a_r<b_{r-1}<b_r=n.
\]
Conversely, pairs with this endpoint order cover all cuts.
With \(b_0=1\), \(a_{r+1}=n\), the cuts from
\(b_{i-1}\) to \(a_{i+1}-1\) are crossed only by pair
\(i\); hence each is necessary. Choosing the \(2r-2\)
interior endpoints determines the matching. This gives
\(\binom{n-2}{2r-2}\), and summing the even binomial
coefficients gives \(2^{n-3}\) for \(n\geq3\).

At a generic point, invert all coordinate differences other
than the pair differences \(u_1,\ldots,u_r\). A cut crossed
only by pair \(i\) has \(F_a=u_i\) times a unit.
Every cut is crossed by some pair, so the localized ideal
is exactly \((u_1,\ldots,u_r)\). These independent linear
forms are a regular sequence. Their local cohomology is
concentrated in degree \(r\), with the nonzero residue class
\(1/(u_1\cdots u_r)\). The shift \([n-1]\) gives
the stated degree.
\end{proof}

These classes occur despite the absence of singular scalar
fields. Cobar records the collision geometry of their
coefficients. A cut crossed by a given equality inverts its
normal coordinate \(u\), so its exceptional restriction
there vanishes. In a normal Koszul model, exterior
multiplication by the normal generator divided by \(u\)
is a contracting homotopy. This explains the same vanishing
directly on the collision complex.

\section{Fields as the differential of a coalgebra}
\label{e1f:sec:open}
\label{e1f:sec:residue-realization}

For a vertex algebra, the pole coefficients supply the
differential of a configuration coalgebra. Its cobar then
has an associative multiplication on the larger coefficient
complexes required by the vacuum calculation.

Take\phantomsection\label{e1f:def:partition-residue-coefficients} \(X=\mathbb A^1\) and a nonzero vertex algebra \(V\)
over \(k\), with current module
\(\mathscr V=\omega_X\otimes_kV\) and the translation
connection of Section~\ref{e1f:sec:pch}. For a partition
\(\pi\) of \([n]\), place one copy of \(\mathscr V\)
at each block. Order the blocks by their largest labels,
called their roots. Extend from distinct block positions
and push along the partial diagonal of \(\pi\). Denote
this right differential module by \(\mathcal V_\pi\).
Its direct image retains all normal derivatives.

For two blocks \(i<j\) in root order, let \(r_{ij}\)
take the principal part of the first state's field at the
second. The other blocks remain as spectators. Taylor
transport of their coordinates is part of this map of
differential modules. The resulting coalgebra coefficient is
\begin{equation}\label{e1f:eq:residue-coalgebra}
 C_n=\bigoplus_\pi\mathcal V_\pi[|\pi|],
 \qquad
 d_C|_{\mathcal V_\pi}
          =\sum_{i<j}(-1)^{i-1}r_{ij},
 \qquad C_0=k.
\end{equation}
A merger reduces the number of blocks by one and has total
degree one. Disjoint mergers cancel with opposite signs.
The remaining terms involve three blocks, whose full
principal parts obey the following identity.

\subsection{Three routes to the same contact module}

Put\phantomsection\label{e1f:def:three-contact-coefficients} \(s=z_1-z_3\), \(t=z_2-z_3\), and \(r=s-t\).
For a vector space \(P\), take the simultaneous coefficient
module
\[
 \mathcal G_3(P)=P[[s,t]][(st(s-t))^{-1}],
\]
where power series are formed before localization. Let
\(E_{s,t}(P)\) be the full contact module spanned by
\(p\,s^{-a}t^{-b}\), \(a,b\geq1\), with its usual
coordinate and normal derivative actions. Write
\(\operatorname{pp}_u\) for the negative powers in a
specified Laurent expansion in \(u\).

The\phantomsection\label{e1f:def:contact-coordinate-transfer} first channel uses contact coordinates \((r,t)\).
To compare it with the other channels one needs the
differential-module map
\[
 J(r^{-a}t^{-b})
   =\sum_{j=0}^{b-1}\binom{a+j-1}{j}s^{-a-j}t^{-b+j}.
\]
It sends \(r^{-1}t^{-1}\) to \(s^{-1}t^{-1}\),
\(\partial_r\) to \(\partial_s\), and the old
\(\partial_t\) to \(\partial_s+\partial_t\).
Applying these derivatives to the simple contact gives the
displayed\phantomsection\label{e1f:def:three-principal-part-routes} formula. Define the three routes, with the first
principal part applied on the right, by
\[
 \begin{aligned}
 L(F)&=J\operatorname{pp}_t\operatorname{pp}_rF(r+t,t)
                                      &&(r\ll t),\\
 R(F)&=\operatorname{pp}_s\operatorname{pp}_tF(s,t)
                                      &&(t\ll s),\\
 N(F)&=\operatorname{pp}_t\operatorname{pp}_sF(s,t)
                                      &&(s\ll t).
 \end{aligned}
\]
They first merge \(12\), \(23\), and \(13\), respectively.
Their targets are now the same contact module, rather than
unidentified spaces in different coordinates.

\begin{lemma}[The three-channel identity]
\label{e1f:lem:three-channel}
For every \(F\in\mathcal G_3(P)\),
\[
 L(F)-R(F)+N(F)=0\quad\text{in }E_{s,t}(P).
\]
\end{lemma}

\begin{proof}
Write \(F=G/(s^at^b(s-t)^c)\). A monomial surviving
two principal parts has total degree at most \(-2\).
Since \(1/(s^at^b(s-t)^c)\) is homogeneous of degree
\(-(a+b+c)\) in every expansion, numerator terms of
degree above \(a+b+c-2\) contribute to none of the three
maps. It suffices to treat polynomial \(G\).

Over \(k[t,t^{-1}]\), the polynomials \(s\) and \(s-t\)
are coprime. Partial fractions reduce the calculation to
polynomials in \(s\), terms \(a_j(t)s^{-j}\), and terms
\(b_j(t)(s-t)^{-j}\), with Laurent coefficients in \(t\)
and the coefficients of \(P\) carried along. A polynomial
in \(s\) contributes nothing. For \(a_j(t)s^{-j}\),
the \(r\)-expansion is regular, so \(L=0\) and
\(R=N=s^{-j}\operatorname{pp}_t(a_j)\).

For \(b_j(t)(s-t)^{-j}\), the expansion in \(s\ll t\)
is regular in \(s\), so \(N=0\). In the other domain,
\[
 (s-t)^{-j}
    =\sum_{n\geq0}\binom{j+n-1}{n}s^{-j-n}t^n.
\]
Only terms \(t^{-\beta}\) of \(b_j\) with \(\beta\geq1\)
contribute. Each gives
\(\sum_{n<\beta}\binom{j+n-1}{n}s^{-j-n}t^{n-\beta}\)
in \(R\), which is exactly \(J(r^{-j}t^{-\beta})\)
in \(L\). This proves the identity on every summand.
\end{proof}

For \(F=1/(s(s-t))\), the values are
\(L=s^{-1}t^{-1}\), \(R=0\), and
\(N=-s^{-1}t^{-1}\), recovering the scalar flags.
The coordinate transfer matters already at the next pole
order: for \(F=1/((s-t)^2t^2)\),
\[
 N=0,\qquad L=R=s^{-2}t^{-2}+2s^{-3}t^{-1}.
\]
Merely renaming \(r\) as \(s\) would miss the second
normal derivative coefficient.

The\phantomsection\label{e1f:def:residue-cut-cobar} identity proves \(d_C^2=0\) in
\eqref{e1f:eq:residue-coalgebra}. On a cut open, a
partition crossing that cut has zero restriction. Every
remaining partition is the external product of a partition
on each side. This gives a strong cut coproduct, commuting
with the differential. Two successive coproducts restrict
to the same open set and are coassociative. The exceptional
collision identifications and the vacuum maps are those
of Theorem~\ref{ors:thm:formal}. Its associative algebra is
\[
 A=\Omega_{\mathrm{cut}}C,
 \qquad \overline A_n\simeq R\Gamma_{Z_n}(C_n)[-1].
\]

\subsection{The outer collision and the full triple product}

At arity three, the finite cut contraction leaves the
two-term coefficient complex
\begin{equation}\label{e1f:eq:residue-ternary}
 G_3=\bigl[\mathcal V_{2|13}[1]
             \xrightarrow{-r_{2,(13)}}\mathcal V_{123}\bigr],
\end{equation}
in degrees \(-1,0\). The first outer residue
\(h=r_{13}:\mathcal V_{1|2|3}\to\mathcal V_{2|13}\)
is therefore a map of degree \(-1\) into \(G_3\).
Writing\phantomsection\label{e1f:def:residue-composite-channels}
\[
 L=r_{(12),3}r_{12},\qquad
 R=r_{1,(23)}r_{23},\qquad
 N=r_{2,(13)}r_{13},
\]
the three-channel identity becomes
\[
 \partial h=-N=L-R.
\]
It is a homotopy on the collision coefficient complex.
It is not a degree-minus-one operation into the original
degree-zero state module.

Nor is it a nullhomotopy of the common triple product.
Let \(\mathcal T=\mathcal V_{1|2|3}\),
\(\mathcal B=\mathcal V_{2|13}\),
\(\mathcal P=\mathcal V_{123}\), and let
\(q:\mathcal P\to G_3\) include degree zero.
Proposition~\ref{ors:prop:derived-triple} identifies the
product of three singleton inputs, under the finite
retract, with \(qL\). If \(L\ne0\), this is a nonzero
derived map.

Indeed, the triangle
\[
 \mathcal B\xrightarrow{-r_{2,(13)}}\mathcal P
       \xrightarrow{q}G_3\longrightarrow\mathcal B[1]
\]
would turn \(qL=0\) into a factorization
\(L=-r_{2,(13)}g\), with
\(g\in H^0\operatorname{RHom}(\mathcal T,\mathcal B)\).
Both coefficients are ordinary modules, so \(g\) would be
an ordinary module map. In coordinates
\(u=z_1-z_3\), \(v=z_2-z_3\), \(s=u-v\), the known
factorization of \(L\) passes through an \(s\)-torsion
module, while this proposed one passes through a
\(u\)-torsion module. The function \(v\) is invertible
before their last collision. Apply
Lemma~\ref{e1f:lem:separation} to \(s,-u\), whose sum
is \(-v\). It would give \(L=0\), a contradiction.
The finite retract sends the fully cut triple product
to \(qL\): its two residue corrections merge \(12\)
and then \((12)3\), and the two suspension signs cancel.

On three vacuum states, the multiplier \(1/(u(u-v))\)
has successive residues \(1/v\) and \(1\) in that
channel. Its output is the nonzero vacuum. The full
derived triple product is consequently nonzero for every
such vertex algebra over \(k\), including the Heisenberg
algebra at level zero
(Corollary~\ref{ors:cor:vacuum-triple}).

The finite contraction extends to every arity. The transferred
\(A_\infty\) operation on singleton inputs is the residue
of the first field at the last, retaining the intervening
fields as spectator blocks. In suspended coordinates it is
\[
 B_n(sx_1,\ldots,sx_n)=s\,r_{1n}(x_1,\ldots,x_n)
\]
(Theorem~\ref{ors:thm:formal}). This describes the
transferred operations on the field complexes; multiplication
in the cobar itself is strictly associative.

The construction also applies over a commutative
\(k\)-algebra of parameters when the state module is flat
and the vacuum spans a free direct summand. It includes
the flat Heisenberg, affine and Virasoro families at their
special levels (Corollary~\ref{ors:cor:vertex}). If fields
are local after a rational unitary exchange, a residue
between nonadjacent blocks first passes the initial field
through the intervening fields. Translation covariance,
Yang--Baxter and the field hexagon identities make these
passages compatible with the residue system
(Theorem~\ref{ors:thm:braided}). The allowed positional
poles are on collision diagonals. Exchanges with other
pole loci or nontrivial double exchange need different
configuration data.

\section{Realizing operations on the collision complexes}

The residue construction gives an algebra and actual
operations on its coefficients. Comparing it with an
independently specified system of fields requires a map
between the corresponding operation complexes.

On\phantomsection\label{e1f:def:diagonal-operation-insertion} the word category, diagonal insertion gives a precise
representation theorem. For \(M\in D(\mathcal D_X)\),
let \(iM\) have coefficients \(\Delta_{n*}M\) and empty
coefficient zero. For any word system \(G\),
Theorem~\ref{oco:thm:operations} identifies
\[
 \begin{aligned}
 &\operatorname{RHom}_{\mathcal W_X}
       (iM_1\star\cdots\star iM_n,G)\\
 &\hspace{8mm}\simeq
 \operatorname{RHom}_{\mathcal D_{X^n}}
       \left(j_{n*}j_n^*(\boxtimes_iM_i),G_n\right).
 \end{aligned}
\]
For \(G=iP\), these are the derived chiral operations,
with their actual compositions. For a field cobar, the
target is instead its coefficient complex \(G_n\).
The representation theorem specifies the maps into that
target; it does not choose which map realizes a proposed
field operation.

To state the further comparison, fix a dg operation system
\(\mathcal Q^{\mathrm{adm}}\) on state modules. Its
singular operations have a finite common pole bound on
each input tuple and compatible expansions, and its domain
is closed under the substitutions under consideration.
Also fix a coefficient category \((\mathcal C,\star,K)\),
either the word category or a specified ordered collision
category. A realization assigns objects \(\mathsf FM\)
and natural maps
\[
 \mathcal Q_n^{\mathrm{adm}}(M_1,\ldots,M_n;P)
 \longrightarrow
 \operatorname{RHom}_{\mathcal C}
       (\mathsf FM_1\star\cdots\star\mathsf FM_n,\mathsf FP).
\]
They must commute with differentials, substitutions, cuts
and the specified collision restrictions. Spectator Taylor
transport and the local vacuum are part of those identities.
On chosen dg models one may ask for injectivity and its
image; the derived question asks for a quasi-isomorphism on
a specified sector. Agreement only in degree zero does not
answer that question.

For binary operations in the weakly translation-equivariant
word category, the degree-zero comparison is known:
translation-covariant fields with finite poles are exactly
the maps \(iM\star iN\to iP\)
(Theorem~\ref{sgc:thm:binary}). The maps themselves must be
equivariant; weak equivariance of their objects alone is
insufficient. This binary identification cannot send full
field composition to consecutive residue composition on
unchanged degree-zero diagonal objects while preserving
the normalized vacuum. The scalar test gives \(1\) and
\(0\) in the two consecutive channels, and
Proposition~\ref{e1f:prop:derived-heart} gives the general
obstruction. The collision complexes retain the third
channel that this attempted comparison would omit.

\subsection{Keeping the runs of an ordered coefficient word}

For a local vertex algebra \(V\) and a unital associative
algebra \(B\) with derivation \(T_B\), the ordered product
on \(V\otimes B\) has coefficients
\[
 m_n^V(v_1,\ldots,v_n;z)\otimes
                  \prod_{i=1}^n e^{(z_i-z_n)T_B}b_i.
\]
The \(B\)-factors stay in input order. The derivation
identity proves every consecutive substitution identity;
their Taylor factors add no poles to those of the fixed
\(V\)-tuple. For dg coefficients the same formula uses
the Koszul shuffle between the \(V\)- and \(B\)-factors.

For a partition \(\pi\), read the name of its block at
each input position. A \emph{run} is a maximal consecutive
string of the same block name. Keep one \(V\)-state per
block and one \(B\)-factor per run. Thus the coefficient
on the stratum of \(\pi\), before localization and
diagonal direct image, is
\[
 \omega_{X^\pi}\otimes V^{\otimes|\pi|}
                           \otimes B^{\otimes\ell(\pi)},
\]
where \(\ell(\pi)\) counts runs. The connection in a
block coordinate includes \(T_V\) on that block and the
sum of \(T_B\) on its runs.

When blocks merge, multiply exactly those coefficient
factors that become one consecutive run. Associativity
makes these maps commute with successive coarsening.
The derivation law makes them commute with the summed
connections. At \(13|2\), the block names are \(a,b,a\),
so the coefficient is \(V^{\otimes2}\otimes B^{\otimes3}\).
The first merger leaves the three \(B\)-factors unchanged;
the final merger multiplies them as \(abc\).

Combining these maps with the \(V\)-field residues gives
the run coalgebra and its strict associative cobar
(Chapter~\ref{sgc:chapter}). Each residue includes the
Taylor exponential on all runs of the translated block.
For a fixed polar input, taking a principal part uses
finitely many Taylor coefficients, so the output remains
in the algebraic tensor coefficient even when \(T_B\)
is not locally finite.

There is a precise reason that these three run factors
cannot always be compressed into two algebraic states.
For the free differential algebra on
\(a_i,b_i,c_i\), translated by increasing the subscript,
consider
\[
 F(t)=(e^{tT_B}a_0)b_0(e^{tT_B}c_0).
\]
An algebraic two-state expression is a finite Laurent-linear
sum of \((e^{tT_B}x)y\) and \(x(e^{tT_B}y)\), with fixed
states \(x,y\). Projecting words \(a_ib_jc_k\) to
\(p^iq^jr^k\) sends \(F(t)\) to \(e^{t(p+r)}\).
Every proposed binary expression becomes a finite sum
\(f(t)P(p,q,r)e^{tL}\), with
\[
 L\in\{0,p,p+q,p+q+r,r,q+r\}.
\]
For \(D=\partial_p+2\partial_q+4\partial_r\), their
exponential eigenvalues are \(st\) with
\(s\in S=\{0,1,3,4,6,7\}\). If \(d\) bounds the
polynomial degrees, \(\prod_{s\in S}(D-st)^{d+1}\)
kills the finite sum and sends the required coefficient to
\[
 (80t^6)^{d+1}e^{t(p+r)}\ne0.
\]
This proves the compression obstruction. It also occurs
in the double-pole channel of \(H_\kappa\otimes B\)
for \(\kappa\ne0\)
(Theorem~\ref{sgc:thm:compression-obstruction}). The
assertion concerns this algebraic two-state framing; it
does not exclude other contextual or completed coefficients.

Naturality imposes a different constraint when later
spectators must remain arbitrary. For \(m\) unfilled
positions, the coefficient context is
\(W_B(m)=B^{\otimes(m+1)}\), evaluated by
\[
 (b_0\otimes\cdots\otimes b_m)(x_1,\ldots,x_m)
                  =b_0x_1b_1\cdots x_mb_m.
\]
A natural quotient retaining every such evaluation must
be injective (Proposition~\ref{sgc:prop:context-forced}).
To see this, adjoin free symbols \(X_1,\ldots,X_m\) to
\(B\). The coefficient of the word
\(X_1\cdots X_m\) is precisely \(B^{\otimes(m+1)}\).
Naturality carries a kernel element to that extension,
and evaluation at the free symbols would kill its injective
word coefficient. Thus the kernel was zero. Truncating
above symbol degree \(m\) proves the same statement with
finite-dimensional extensions when \(B\) is finite-dimensional.
A quotient at a single fixed algebra can be smaller;
the naturality requirement is essential.

\subsection{A merger using the given middle state}

If\phantomsection\label{e1f:def:middle-state-exchange} \(T_B=0\), a specified exchange provides another
coefficient construction. For an ordinary nonzero unital
algebra \(B\), define
\[
 c_B(a\otimes b)=ab\otimes1+1\otimes ab-a\otimes b.
\]
Its unit identities are
\(c_B(a\otimes1)=1\otimes a\) and
\(c_B(1\otimes a)=a\otimes1\). Applying them twice gives
\(c_B^2=1\). Expansion using associativity proves the
braid relation and the two multiplication compatibilities
of Lemma~\ref{ctx:lem:braid}. If \(\tau_B\) is the flip,
the field exchange \(\mathcal S_B=c_B\tau_B\) obeys
\[
 (\mathcal S_B)_{21}\mathcal S_B=1,
 \qquad \mu_B\mathcal S_B=\mu_B\tau_B,
\]
and fixes each vacuum factor. The multiplication
compatibilities are its field hexagons.

Give \(V\otimes B\) the fields \eqref{e1f:eq:VB} and
let this exchange act on the \(B\)-factors. Locality of
\(V\), the last displayed equality and associativity in
\(B\) prove the exchanged field identities. Theorem~\ref{ors:thm:braided}
therefore gives their residue coalgebra with one state per
occupied position.
For the order \((13,2)\), the coefficient merger is
\begin{equation}\label{e1f:eq:contextual-outer}
 \begin{aligned}
 q_{13|2}:B^{\otimes3}&\longrightarrow B^{\otimes2},\\
 q_{13|2}(a,b,c)&=abc\otimes1+a\otimes bc-ab\otimes c.
 \end{aligned}
\end{equation}
Its output product is \(abc\). For
\(a=c=E_{12}\), \(b=E_{21}\), the value is
\[
 E_{12}\otimes1+E_{12}\otimes E_{22}
                         -E_{11}\otimes E_{12},
\]
whose product is \(E_{12}\), even though \(ac=0\).
Changing to root order \((2,13)\) applies \(c_B\) and gives
\[
 c_Bq_{13|2}(a,b,c)=ab\otimes c+1\otimes abc-a\otimes bc
\]
(Proposition~\ref{ctx:prop:qfunctor}). The inserted
\(1_B\)'s belong to retained positions; they are not
empty-configuration units.

This map uses the actual middle state. It therefore poses
a different problem from preserving every evaluation on
an unfilled context. The run coefficient has three
\(B\)-factors; the exchanged coefficient has two; their
fully multiplied value is \(abc\). Equality of that
value does not identify the two partial-stratum objects.
The assumption \(T_B=0\) also leaves the free differential
compression obstruction intact.

\subsection{Homotopies relative to the field comparison}

For\phantomsection\label{e1f:def:relative-residue-comparison} the Heisenberg--matrix run construction, let \(F\)
be the diagonal input system with its binary residue
\(r:F\star F\to F\), and let \(A\) be the run cobar
with multiplication \(\mu\). The one-block inclusion
\(\lambda:F\to A\) is a chain map. If \(\tau\) denotes
the degree-one map from the coalgebra to its cobar
generators, the interval two-block coefficient gives
\(f_2:F^{\star2}\to A\) of degree \(-1\), with
\[
 \partial f_2=\mu(\lambda,\lambda)-\lambda r.
\]
The outer contact \(\mathsf R_{13}\) gives
\(h=-\tau\mathsf R_{13}:F^{\star3}\to A\), with
\[
 \partial h=\lambda\mathfrak a,
 \qquad \mathfrak a=r(r,1)-r(1,r)=L-R=-N.
\]
The sign includes the root transfer of the contact module;
it is fixed on that module, rather than by renaming a
polar variable (Chapter~\ref{ncc:chapter}).

The four-input identity must include the binary comparison
homotopy. Its degree-minus-one boundary is
\[
 \begin{aligned}
 \mathcal O_4={}&h(r,1,1)-h(1,r,1)+h(1,1,r)\\
   &-\mu(h,\lambda)-\mu(\lambda,h)
       +f_2(\mathfrak a,1)+f_2(1,\mathfrak a).
 \end{aligned}
\]
Every term maps \(F^{\star4}\) to \(A\).
The five binary trees cancel in the associator identity;
using the displayed equation for \(f_2\) then gives
\(\partial\mathcal O_4=0\). Without the two \(f_2\)
terms the expression need not even be closed.

There is an explicit degree-minus-two nullhomotopy. Let
\([i,j]\) record the unsigned first contact and its run
map, and let angle brackets indicate the coefficient on
the displayed partition. Then
\[
 K_4=\tau\bigl(
       -\langle[1,3]\mid2\mid4\rangle
       +\langle[1,4]\mid2\mid3\rangle
       +\langle1\mid[2,4]\mid3\rangle\bigr)
\]
satisfies \(\partial K_4=\mathcal O_4\), including its
quadratic cobar terms (Chapter~\ref{ncc:chapter}).
Closedness of the three-channel expression alone would
not construct this map.

These are comparisons from powers of the diagonal input
to a larger cobar. They give neither an \(A_\infty\)
structure on the unchanged input nor higher multiplication
in place of the strict multiplication of \(A\). Transfer
to the smaller field complexes is a separate
\(A_\infty\) construction, as in
\eqref{e1f:eq:residue-ternary}.

\subsection{Higher operations on prescribed generators}

An\phantomsection\label{e1f:def:prescribed-generator-diagram} algebraic resolution can require higher terms for the
same reason that the collision comparison retains more
than its binary values. Let \(H\) be a dg coefficient
algebra of weight zero, and consider a diagram
\((T_HV,d)\) augmented to \(H\). The differential preserves
weight, and the diagram maps are induced on the generators.
Its linear part \(d_1\) makes \(V\) a complex of
\(H\)-bimodules. Assume positive integer weights and
\(K\)-flatness on both sides: tensoring with an acyclic
module on the opposite side must remain acyclic.

Theorem~\ref{cgr:thm:resolution} resolves this prescribed
generator complex by a semi-free diagram \(P\to V\)
and lifts the differential to a quasi-isomorphism
\[
 (T_HP,D)\longrightarrow(T_HV,d).
\]
Semi-free generators are attached in an order whose
differentials use earlier generators. A decomposable value
of weight \(w\) uses lower-weight generators and has
length at most \(w\). The lift may consequently have
higher word-length terms while requiring no completion.
Applying this result to chiral data still requires the
specified geometric diagram and tensor-flatness at its
collisions and coefficient changes.

The\phantomsection\label{e1f:def:higher-generator-differential} zero-parameter example of Proposition~\ref{cgr:prop:example}
already shows the higher terms over \(H=k\). On free
generators with \(|x|=|v|=0\), \(|u|=|y|=-1\),
\(|z|=-2\), put
\[
 Dx=Dv=0,\qquad Du=v,\qquad Dy=x^2+vx,
 \qquad Dz=xy-yx-xux+ux^2.
\]
The differential of the quadratic part of \(Dz\) is
\(xvx-vx^2\); the cubic part has the opposite differential.
Hence \(D^2z=0\). The change \(Y=y-ux\) gives
\(DY=x^2\), \(Dz=xY-Yx\). The two presentations describe
the same derived algebra, but the prescribed larger
generator complex needs the cubic correction.

The transverse sphere provides a different example.
In its finite angular-mode dg algebra, the triple Massey
product of two holomorphic coordinates with a degree-one
mode is \(-1\) modulo the coordinate ideal
(Theorem~\ref{rbt:thm:nonformal}). Its nonzero constant
class forces higher products in the effective coefficient
model. To realize them on a chiral defect one must still
construct their localized inputs, collision supports and
comparison maps (Section~\ref{rbt:sec:chiral}).

\section{Recovering the symmetric theory}
\label{e1f:sec:symmetric}

A symmetric chiral algebra supplies more collision maps
than a word system requires. Its passage to the word
category is most transparent on its coalgebra.
Francis--Gaitsgory Chevalley chains give an equivalence
from Lie objects to cocommutative coalgebras for the chiral
tensor on Ran space \cite{chhop:FG}*{Theorem~5.1.1}.
The coalgebra factorizes exactly when the Lie object is
supported on the singleton copy of \(X\)
\cite{chhop:FG}*{Theorem~5.2.1}. Those Lie objects are
chiral Lie algebras without a specified vacuum.

Restrict their collision maps to monotone surjections and
project their all-subset coproduct onto consecutive cuts.
Coassociativity survives: either order of two projections
retains the same three consecutive blocks. Factorization
also survives on those cut opens. The cut cobar of the
restricted coalgebra is therefore an associative chiral
algebra in Definition~\ref{e1f:def:e1}. This constructs a
word algebra from the chiral Lie algebra. A vacuum framing
additionally compares its operations with the original
simultaneous fields.

\subsection{Why restriction must precede cobar}

Restricting\phantomsection\label{e1f:def:free-commutative-singleton-species} an algebra before taking bar need not give the
same coefficient complex as restricting its bar. On the
separated coefficient species, take the free commutative
algebra on one singleton generator. It has \(A(I)=k\)
for every finite set \(I\), with identity products.
For \(1<2\), its cut bar has the two-term complex
\[
 k[1|2]\longrightarrow k[12],
\]
whose map is an isomorphism. The all-subset bar instead has
\[
 k[1|2]\oplus k[2|1]\longrightarrow k[12],
\]
with both words mapped to \([12]\). It retains
\([1|2]-[2|1]\), the Chevalley class of the two commuting
inputs. Dropping the reverse word is not a chain map:
it kills \([2|1]\) but not its differential.

Restriction of the coalgebra keeps this component complex.
Taking its cut cobar and then its cut bar recovers it by
bar--cobar duality. It therefore preserves the class that
the bar of the restricted algebra loses.

In the other direction, symmetry on distinct points alone
does not restore nonconsecutive collision maps.

\begin{proposition}\label{e1f:prop:descent}
There is a cut coalgebra with ordinary symmetry on every
fully distinct locus which has no extension compatible
with all surjections.
\end{proposition}

\begin{proof}
Choose a differential module \(M\) with
\(j^*(M\boxtimes M)\ne0\). Let \(D(M)\) be its
diagonal word system, zero at the empty word, and take
\[
 C=K\oplus\bigoplus_{r\geq1}D(M)[1]^{\star r}
\]
with deconcatenation. On a fully distinct locus only
singleton blocks remain, and their factors have the
ordinary Koszul symmetry.

The nonmonotone surjection with word \(121\) has diagonal
\(\delta(x,y)=(x,y,x)\). Every summand of \(C_3\) with
at least two consecutive blocks inverts a difference that
vanishes under \(\delta\). Its exceptional pullback is
therefore zero. The surviving one-block summand of
\(\delta^!C_3\) is supported on \(x=y\).
But \(C_2\) also contains the nonzero open coefficient
\(j_*j^*(M[1]\boxtimes M[1])\). These two objects cannot
be the coefficients identified by the surjection.
\end{proof}

A reverse comparison thus needs an extension of the whole
bar coalgebra along every surjection, with a compatible
cocommutative coproduct and factorization on disjoint
loci. Under those hypotheses the Francis--Gaitsgory
theorems recover a chiral Lie algebra. Descent of just
the underlying modules leaves the required coproduct
unspecified, and averaging is generally insufficient.

\begin{theorem}\label{e1f:thm:average}
Let \(C\) be a finite-dimensional ordinary coassociative
coalgebra, with dual algebra \(A\). The coproduct
\(\tfrac12(\Delta+\tau\Delta)\) is coassociative exactly
when \([A,[A,A]]=0\).
\end{theorem}

\begin{proof}
The dual multiplication is \(a\circ b=(ab+ba)/2\).
Its associator is
\[
 (a\circ b)\circ c-a\circ(b\circ c)
 =\frac14(bac+cab-acb-bca)=\frac14[b,[a,c]].
\]
Thus it vanishes for all inputs exactly under the stated
condition. In \(\operatorname{Mat}_2(k)\), the choice
\(a=E_{11}\), \(b=E_{12}\), \(c=E_{21}\) gives
\((E_{22}-E_{11})/4\), so averaging fails there.
\end{proof}

\begin{example}\label{e1f:ex:noncommutative-average}
The condition allows noncommutative algebras. In
\(\operatorname{Mat}_3(k)\), take
\[
 A=kI_3\oplus kE_{12}\oplus kE_{23}\oplus kE_{13}.
\]
Here \(E_{12}E_{23}=E_{13}\) and
\(E_{23}E_{12}=0\), but \(E_{13}\) commutes with all
of \(A\). Hence \([A,A]=kE_{13}\) and
\([A,[A,A]]=0\). The averaged coproduct of its dual
coalgebra is coassociative despite the original algebra's
noncommutativity.
\end{example}

\subsection{The Chevalley comparison and its normalization}

A\phantomsection\label{e1f:def:chevalley-comparison-input} Lie map gives a structured comparison without averaging
the associative coproduct. Let \(A\) be an augmented dg
algebra in a symmetric coefficient category, let
\(\overline A^-\) be the commutator Lie algebra of its
augmentation ideal, and take a dg Lie map
\(f:\mathfrak g\to\overline A^-\). On suspended inputs,
put
\[
 sx_1\cdots sx_r\longmapsto
 \sum_{\sigma\in S_r}\varepsilon(\sigma;sx_1,\ldots,sx_r)
 [f(x_{\sigma(1)})|\cdots|f(x_{\sigma(r)})].
\]
This maps Chevalley chains, with unshuffle coproduct, to
the tensor bar, with deconcatenation. Splitting the
permutation sum is the same as choosing a subset and
independently ordering it and its complement. Hence it
preserves coproducts. In the differential, the two orders
of each adjacent pair combine into the suspended Lie
bracket, so it is a chain map.

The sum is unnormalized. For degree-zero \(x,y\), its
binary value is \([f(x)|f(y)]-[f(y)|f(x)]\), whose
differential is the bracket image. Dividing this component
by \(2\) while leaving the linear component fixed would
break that equality; dividing by \(r!\) would also change
the unshuffle/deconcatenation comparison.
Theorem~\ref{ogd:thm:CE} constructs this map on the symmetric
monoidal ordered collision category and proves that
\(\operatorname{CE}_\star(L)\to B(U_\star L)\) is an
equivalence for positive Lie objects. Its PBW proof uses
a finite configuration filtration. The PBW symmetrization
does have a factorial normalization; the coalgebra map
just displayed does not.

The commutator and this comparison use symmetry of the
coefficient tensor. The nonsymmetric cut convolution
does not supply it. Nor does the envelope theorem
identify the bar of an arbitrary associative algebra
with the Chevalley chains of its commutator Lie algebra.

\section{Coproducts, the second tensor and centres}
\label{e1f:sec:map}

An associative product combines configurations. A coproduct
of field algebras instead compares states occupying the
same position. Compatibility between them needs a second
tensor and an interchange map with its associativity and
unit identities. Associative bar--cobar duality for one
tensor does not construct this additional tensor.

For an OPE bialgebra
(Definition~\ref{cqe:def:bialgebra}), the same-position
coproduct is a map
\(G\to G\otimes_{\mathcal O_X}G\) preserving the vacuum
and every simultaneous product. It has a counit to
\(\mathcal O_X\). Multiplication of the scalar pole
functions defines the tensor product of OPE algebras.
At a fixed nonzero Heisenberg level, the counit is already
excluded by a singular coefficient.

\begin{proposition}\label{e1f:prop:counit}
If \(\kappa\ne0\), there is no unital vertex-algebra map
\(H_\kappa\to k\).
\end{proposition}

\begin{proof}
Such a map \(\varepsilon\) would preserve the first
singular product and give
\[
 \kappa=\varepsilon(h_{(1)}h)
       =\varepsilon(h)_{(1)}\varepsilon(h)=0,
\]
since the scalar target has no singular products.
\end{proof}

The\phantomsection\label{e1f:def:universal-central-current} universal central state has a coproduct even though
its fixed nonzero specialization has no counit. Let
\(\mathcal H\) be universally generated by \(h\) and
a central state \(C\), with \(TC=0\),
\(h_{(1)}h=C\), and \(h_{(n)}h=0\) for \(n=0\)
and \(n\geq2\). Then
\(\mathcal H/(C-\kappa\mathbf1)=H_\kappa\).
The assignments
\[
 h\longmapsto h\otimes\mathbf1+\mathbf1\otimes h,
 \qquad
 C\longmapsto C\otimes\mathbf1+\mathbf1\otimes C
\]
preserve the generating relations: mixed singular products
between the two tensor factors vanish. Universality gives
a vertex-algebra coproduct. Both triple iterations send
each generator to the sum of its three copies, proving
coassociativity. Sending \(h,C\) to zero gives the counit.
Specializing the two output central states gives
\[
 H_{\kappa_1+\kappa_2}\longrightarrow
                         H_{\kappa_1}\otimes H_{\kappa_2}.
\]
Thus the coproduct adds levels. It does not preserve a
fixed nonzero level in both output factors.

\subsection{The tensor on collision coalgebras}

On the ordered collision category there is a diagonal tensor
\(\otimes^!\) on convolution coalgebras, with dualizing unit
\(J\). Its coefficient at one point is \(\omega_X[1]\),
and its empty coefficient is \(k\). Bar--cobar transports
this tensor to the product
\[
 A\diamond A'=\Omega(BA\otimes^!BA'),
 \qquad \mathbf1_\diamond=\Omega J.
\]
This differs from convolution, whose unit is \(K\).
A monoid for \(\otimes^!\) among convolution coalgebras
is called a \emph{transverse monoid}. Its multiplication
must preserve the cut coproduct through the geometric
interchange; it is not a second copy of cut multiplication.

The ordinary diagonal tensor need not agree with the
tensor of fields carrying a cross exchange. At two
separated positions, the two coproduct components of a
factorizing coalgebra determine an exchange by inverting
one component and composing with the other. A coalgebra
map with prescribed singleton map must intertwine these
exchanges. Proposition~\ref{cqb:prop:diagonal-obstruction}
gives constant unitary spectral forms for which the field
tensor and the ordinary diagonal tensor have different
exchanges on the same state tensors. Hence no comparison
inducing their singleton identity exists in that example.

There is a positive comparison when the tensor retains
the spectral regrouping. On the admissible covariant word-coefficient
domain of Definition~\ref{stp:def:covariant}, the state
coactions supply rational differential-module passage
maps, with finite pole order on each state, compatibility
with collision pullbacks, and continuity for any chosen
completion. For a unitary spectral form, insert the
inverse regrouping in the diagonal tensor's cut
coproduct. This gives the monoidal tensor
\(\otimes^!_\sigma\), with unit \(J\)
(Proposition~\ref{stp:prop:monoidal}). Its coefficient
objects are those of the diagonal tensor; its coproduct
has the additional passage map.

For the covariant field tensor \(V_1\boxtimes_\rho V_2\),
the comparison is
\[
 C(V_1\boxtimes_\rho V_2)
       \xrightarrow{\ \sim\ }
 C(V_1)\otimes^!_\sigma C(V_2)
\]
(Theorem~\ref{tpd:thm:twisted}). It is constructed on
regular lattices and then passed through their comparison
with the residue complexes. It is natural, monoidal,
the identity on singleton states, and preserves the
vacuum marking under \(J\). Outside this coefficient
domain, admissible passage maps and independence of the
presentation remain separate requirements
(Question~\ref{cqb:qn:second-tensor}). A bare coaction
on an arbitrary completed module does not supply them.

\subsection{The universal property depends on the action tensor}

For symmetric convolution on the ordered collision
schemes, an associative algebra \(A\) has a convolution
centre \(Z^*(A)\). It is the final \(E_2\)-algebra acting
through maps \(F\star A\to A\) that preserve the algebra
structure (Theorem~\ref{oco:thm:center}). It represents
the internal endomorphisms of the bimodule \(A\), and
its unit is a map from the empty-configuration object
\(K\). The centre itself retains configuration
coefficients.

On the nonsymmetric word category, the corresponding
scalar centre is
\(R\operatorname{Hom}_{A\text{-}A}(A,A)\). It is final
among actions of \(E_2\)-algebras in \(D(k)\), acting
by \(F\otimes_k A\)
(Theorem~\ref{oco:thm:word-center}). This is a different
coefficient and action domain from the preceding internal
centre.

An interior field can act on a boundary field at the
same position of the curve. Its geometric action uses
the diagonal tensor on coalgebras, or \(\diamond\)
after cobar. For a coalgebra \(C\) in the symmetric
convolution category, there is an
ambient universal transverse monoid \(\mathcal E(C)\)
acting by \(F\otimes^!C\to C\), with no factorization
requirement on the actor
(Theorem~\ref{tvp:thm:universal}). Requiring a final
\emph{factorizing} actor changes the existence question.

Such an actor need not exist for a general factorizing
coalgebra. The split Chevalley coalgebra \(C(W)\) of
an abelian diagonal Lie object, with finite-dimensional
singleton \(W\otimes\omega_X[1]\), gives a counterexample
when \(\dim W\geq2\)
(Theorem~\ref{scn:thm:no-centre}). Each decomposition
of \(W\) into two summands gives a factorizing action
through its complementary projectors. On \(W=k^2\),
take
\[
 P=\begin{pmatrix}1&0\\0&0\end{pmatrix},
 \qquad Q=\begin{pmatrix}1&1\\0&0\end{pmatrix}.
\]
Both are idempotent, but \(PQe_2=e_1\) and \(QPe_2=0\).
A common factorizing recipient would make their images
commute on the collision diagonal, contradicting these
values. This rules out a final recipient for all such
actions.

The split coalgebra has no nonzero local vacuum. It
therefore does not settle existence for coalgebras of
unital flat chiral algebras. On that class the ordinary
chiral centre is final among factorizing actors bounded
above by the heart (Theorem~\ref{htd:thm:below}); the
unrestricted derived problem remains
Question~\ref{htd:qn:centralizer}. Two cases of it are decided
negatively. For the algebraic Heisenberg coalgebra at nonzero
level on the affine line over a field of characteristic zero,
the pointed-action functor is not representable
(Theorem~\ref{pso:thm:sphere}). For an ordinary unital vertex
algebra over \(\mathbb C\) on the affine line, with algebraic
states in degree zero and a nonzero vertex derivation, the
category of coherent factorizing transverse monoid actions on
its Chevalley--Cousin coalgebra has no final object
(Theorem~\ref{env:thm:ordinary}). The unrestricted form
of Question~\ref{tvp:qn:local} must be distinguished
from this remaining vacuum-bearing case. A mixed
bulk--boundary realization also requires the comparison
of Conjecture~\ref{intro:conj:ht-centre}; neither
convolution finality nor the ambient transverse monoid
constructs it.

The carriers now give the following comparisons. In
each row, an equivalence includes only the domain and
structures stated there.
\medskip
\begingroup
\renewcommand{\arraystretch}{1.2}
\begin{longtable}{
 >{\raggedright\arraybackslash}p{.29\textwidth}
 >{\raggedright\arraybackslash}p{.31\textwidth}
 >{\raggedright\arraybackslash}p{.31\textwidth}}
\toprule
Input & Output & Comparison\\\midrule
\endfirsthead
\toprule
Input & Output & Comparison\\\midrule
\endhead
\bottomrule
\endfoot
Strict associative chiral map with normalized vacuum
 & Zero module & Theorem~\ref{e1f:thm:vacuum}\\
Associative object in vertex algebras
 & Regular commutative differential algebra
 & Equivalence, Theorem~\ref{e1f:thm:va}\\
Chiral algebra with vacuum and translation on \(\mathbb A^1\)
 & Vertex algebra & Equivalence, \cite{chhop:BD}*{\S0.15}\\
Regular ordered product system
 & Associative differential algebra
 & Equivalence, Proposition~\ref{e1f:prop:regular}\\
Ordered product system
 & System with a finite exchange
 & Extra data; obstruction in Proposition~\ref{e1f:prop:noexchange}\\
Chiral Lie algebra without vacuum
 & Associative cut algebra
 & Chevalley chains, restriction, then cobar\\
Associative cut algebra
 & Chiral Lie algebra
 & Needs a cocommutative factorizing extension of its bar\\
Coassociative coalgebra
 & Cocommutative coalgebra by averaging
 & Criterion of Theorem~\ref{e1f:thm:average}\\
Augmented collision algebra, empty coefficient \(k\)
 & Coaugmented collision coalgebra, empty coefficient \(k\)
 & Bar--cobar equivalence; it restricts to factorizing bars\\
Vertex fields or rational unitary exchanged fields
 & Residue coalgebra and associative cobar
 & Theorem~\ref{ors:thm:formal}; framing uses the specified operations\\
\end{longtable}
\endgroup
\medskip

Bar--cobar equivalence also applies outside the factorizing
sector. For example, put a primitive coproduct on
\(K\oplus D(\omega_X)\), with positive coefficients on
the small diagonals. At two distinct positions the
coalgebra vanishes but the external product of its
singleton coefficients does not. Factorization fails.
The equivalence restricts to factorizing bars precisely
because \(B\Omega C\simeq C\) preserves the actual
coproduct maps.

\section{Independent inputs and normal collision jets}

A field product can identify distinct input tensors.
Its geometric coefficients must specify whether they
retain those tensors or only their evaluated product.
Let\phantomsection\label{e1f:def:independent-input-regular-lattice} \(V\) be a simultaneous ordered product system on
\(X=\mathbb A^1\), with a finite common denominator
for each ordered tuple. Keep the original input module
\[
 \mathcal H_I=\mathcal O_{X^I}(*D_I)\otimes_k V^{\otimes I},
 \qquad D_I=\bigcup_{i<j}\{z_i=z_j\}.
\]
For an interval partition \(\pi\), let \(M_\pi\)
apply the simultaneous product within each block and
retain the different block outputs. Expand in their
internal separations with differences between the block
roots inverted. Define
\[
 \mathcal K_I=\bigcap_{\pi\ \mathrm{interval}}
 M_\pi^{-1}\bigl(\text{series regular in all internal
                         separations of }\pi\bigr).
\]
This is a lattice inside the prescribed independent
inputs. Here ``lattice'' means a coordinate submodule
whose localization recovers those inputs; finite
generation is not required.

For a constant unital algebra \(B\), at two positions
and with the common coordinate suppressed, multiplication
has the section \(b\mapsto1\otimes b\). Thus
\[
 \mathcal K_2
 =\ker(B\otimes B\to B)[u,u^{-1}]
                      \oplus(1\otimes B)[u],
 \qquad u=z_1-z_2.
\]
The kernel directions retain arbitrary poles because
their field product is zero. Passing to the image of
multiplication would erase these directions before
the collision is considered.

The order matters for the same reason. For
\(B=k[x]/(x^2)\), the input
\[
 s=\frac{x\otimes x\otimes1}{z_1-z_3}
\]
lies in \(\mathcal K_3^{123}\). Every interval block
containing both labels \(1,3\) contains label \(2\),
so its product contains \(x^2=0\); in every other
test the denominator is between blocks. In order
\(132\), however, the block \(13\) has product
\(x\), leaving the nonzero two-block coefficient
\(x\otimes x/(z_1-z_3)\) with an internal pole.
Thus \(s\notin\mathcal K_3^{132}\). Comparing the
lattices needs a transition map, rather than equality
under the identity on inputs.

On consecutive strata their collision maps are exact.
For an adjacent merger with normal coordinate \(u\),
the map \(c_*\) takes the constant internal coefficient
of the corresponding regular partial product. It fits
into
\[
 0\longrightarrow\mathcal K_I
    \xrightarrow{\ u\ }\mathcal K_I
    \xrightarrow{\ c_*\ }\mathcal K_J
    \longrightarrow0.
\]
Insertion of the vacuum splits the last map on
coefficient complexes. This sequence is linear for
coordinate functions and compatible with tangent
translations. Multiplication by \(u\) is not a full
differential-module map, since \([\partial_u,u]=1\).
Iterated normal Koszul resolutions give the exceptional
collision comparisons after the density normalization
\((\omega_{X^I}\otimes\mathcal K_I)[|I|]\).
Localization at a consecutive cut gives the external
product of its two lattices. Chapter~\ref{cl:chapter}
proves these collision and factorization statements,
including their compatibility.

Nonconsecutive coefficients must also retain spectator
poles and normal derivatives. At Heisenberg level
\(\lambda\), the two contractions in
\(:h^2:(z_1)h(z_2)h(z_3)\) have vacuum coefficient
\[
 \frac{2\lambda^2}{(z_1-z_2)^2(z_1-z_3)^2}.
\]
Set \(u=z_1-z_3\), \(v=z_2-z_3\), and expand with
the middle point held apart. Its outer principal part is
\[
 \operatorname{pp}_u
       \frac{2\lambda^2}{u^2(u-v)^2}
   =2\lambda^2\left(v^{-2}u^{-2}+2v^{-3}u^{-1}\right).
\]
The double contact includes a normal derivative; the
simple contact has a spectator pole of order three.
Keeping only a residue with the original order-two
spectator bound would miss part of this coefficient.

\section{Common pole bounds and continuous substitution}

The constructions above began with simultaneous products.
Weak associativity compares iterated fields after
multiplication by a suitable regularizing power, which
can depend on the chosen states. Its implication for
a common multivariable target is a further question.

\begin{question}\label{e1f:qn:simultaneity}
Does weak associativity on every triple of states in a
nonlocal vertex algebra force every four-state product
to be the expansion of a common germ with poles only
on collision divisors?
\end{question}

A counterexample must be an actual nonlocal vertex
algebra: its fields, vacuum and translation must obey
weak associativity on all states. It must then have
a four-state product with no finite common collision
denominator. Single generation is not required.
The scalar expression
\[
 \exp\left(\frac{(x-y)(y-z)z}{x}\right)
\]
illustrates the distinction between coefficientwise
and common bounds. Its coefficients after expansion
in \(y,z\) are meromorphic in \(x\), but for generic
fixed \(y,z\) it has an essential singularity at
\(x=0\). No finite collision denominator removes it.
This expression by itself supplies none of the field
algebra axioms required for a counterexample.
An associativity bound depending only on the outer
pair of states and uniform in the middle state is
a stronger requirement.

The same issue arises when operations are substituted.
Composing \(\sum_{n\geq0}u^ne_n\) with
\(e_n\mapsto v^{-n^2}w\) gives
\(\sum_{n\geq0}u^nv^{-n^2}w\). Each coefficient has
a finite pole, but no single power of \(v\) clears
all coefficients. There can also be a nonmeromorphic
composite with only linearly growing intermediate
poles: Proposition~\ref{ecd:prop:log} constructs
\(\log(1-v/u)\) from such operations. Thus finite
pole order on each intermediate state does not make
an operation domain closed under substitution.

\subsection{Ordered tensor topology}
\label{e1c:sec:topology}

The\phantomsection\label{e1c:def:ordered-tensor-completion} modes of fields are operators, and their
composition is associative. A topology specifies
which infinite operator expressions this permits.
For complete separated linearly topologized vector
spaces \(E_1,\ldots,E_n\), declare a subspace
\(U\subset E_1\otimes\cdots\otimes E_n\) open when,
for every \(i\) and every fixed suffix
\(e_{i+1},\ldots,e_n\), there is an open
\(P\subset E_i\) with
\[
 E_1\otimes\cdots\otimes E_{i-1}\otimes P
              \otimes e_{i+1}\otimes\cdots\otimes e_n
                       \subset U.
\]
Denote its separated completion by
\(E_1\otimes^{\to}\cdots\otimes^{\to}E_n\).
These tensors are associative with discrete unit
\(k\) \cite{chhop:BD}*{\S3.6.1}.
Their order specifies the corresponding iterated
Laurent-series space when the factors are \(k((t_i))\).

\begin{proposition}\label{e1c:prop:operators}
For a discrete vector space \(W\), give
\(\mathcal E=\operatorname{End}_k(W)\) the neighborhoods
\[
 U_F=\{a\in\mathcal E:a|_F=0\},
                      \qquad \dim F<\infty.
\]
Composition makes \(\mathcal E\) a complete unital
associative algebra for \(\otimes^{\to}\).
\end{proposition}

\begin{proof}
A Cauchy net is eventually constant on each vector
of \(W\). Its values define a unique linear operator,
which proves completeness and separation. For
\(m(a,b)=ab\), the inverse image of \(U_F\) contains
\(\mathcal E\otimes U_F\), and for a fixed \(b\)
it contains \(U_{b(F)}\otimes b\). These are the
two ordered continuity conditions, so composition
extends to the completion. Its two associative
composites agree on the dense algebraic triple
tensor and therefore agree on its completion.
The identity operator gives the unit.
\end{proof}

Every mode on \(W\) belongs to this operator algebra.
This fact imposes no uniform Laurent bound on all
states of \(W\). Substitution of a formal sum of
states is a separate continuity question.

\begin{proposition}\label{e1c:prop:completion}
Let \(\kappa\ne0\), give \(H_\kappa\) the discrete
topology and \(H_\kappa[[u]]\) the \(u\)-adic topology.
There is no continuous substitution on all of
\(H_\kappa[[u]]\) agreeing on finite sums with
\(a(u)\mapsto Y(a(u),u)h\), into a target whose
Laurent-coefficient maps to discrete \(H_\kappa\)
are continuous.
\end{proposition}

\begin{proof}
The partial sums of
\[
 A(u)=\sum_{r\geq0}\frac{(-1)^ru^r}{(r+1)!}T^rh
\]
converge \(u\)-adically. Differentiating the
Heisenberg double pole gives
\[
 Y(T^rh,u)h=(-1)^r(r+1)!\kappa u^{-r-2}\mathbf1
                                  +\text{regular terms}.
\]
Every term of \(A(u)\) therefore contributes
\(\kappa\mathbf1\) to the coefficient of \(u^{-2}\)
after substitution. The first \(N\) terms contribute
\(N\kappa\mathbf1\), which is not eventually
constant. Continuity of substitution and coefficient
extraction would make this a convergent sequence in
the discrete coefficient space, a contradiction.
\end{proof}

Completing\phantomsection\label{e1c:def:formal-diagonal-output} a geometric output instead specifies a
different operation. Let \(G\) be a left differential module,
let \(\mathcal I_\Delta\) be the ideal of a small diagonal,
and choose an output coordinate. The underlying formal
coefficient module is
\[
 \widehat\Delta_{I,*}G
 =\varprojlim_N
       (\mathcal O_{X^I}/\mathcal I_\Delta^N)
                       \otimes_{\mathcal O_X}G.
\]
Here tensoring is performed before the inverse limit.
It allows infinitely many independent state coefficients;
Taylor transport compares choices of output coordinate.
Localizing afterwards at the collision differences
allows a finite denominator for each input tuple.
It does not permit substitution of every completed
state sum. The preceding proposition separates these
two completion requirements.

\subsection{A dual algebra and its geometric comparison}

The first associative dual is the coalgebra
\(B_{\mathrm{cut}}A\). A specified coefficient duality
can transpose its coproduct into an algebra in the
dual coefficient category
(Theorem~\ref{full-thm:koszul}). This transposes the
operations it is defined on. A realization as a
second field system still needs its collision maps,
vacuum and coefficient topology.
A boundary realization also requires an independently
specified module and its action.

An\phantomsection\label{e1f:def:derived-endomorphism-field} algebraic endomorphism model makes the choice of
recipient explicit. For an augmented dg algebra
\(B\), choose a homotopically projective resolution
\(P\to k\): for every acyclic \(B\)-module \(N\),
\(\operatorname{Hom}_B(P,N)\) is acyclic. Set
\[
 E=\operatorname{Hom}_B(P,P),\qquad
 d_Ef=d_Pf-(-1)^{|f|}fd_P,\qquad fg=f\circ g.
\]
This is a unital associative dg algebra.
Applying \(\operatorname{Hom}_B(P,-)\) to \(P\to k\)
identifies its underlying complex with
\(R\operatorname{Hom}_B(k,k)\). Its regular ordered
field is \(Y(f,z)g=fg\), with translation zero and
vacuum \(1_P\). Thus it gives an algebra-valued dual
on the chosen resolution. It has not yet supplied
a comparison of singular chiral operations.

Transferring \(E\) to \(H^*(E)\) retains higher
\(A_\infty\) operations and a quasi-isomorphism;
the binary cohomology product alone need not determine
\(E\). Recovering \(B\) from a double dual also requires
an augmentation-biduality comparison in the chosen
coefficient category.

Field\phantomsection\label{e1f:def:exact-vacuum-contraction} comparisons must also respect the cohomological
degree of the vacuum. In a unital dg field algebra,
suppose \(\mathbf1=dv\) with \(|v|=-1\). Since
\(d\) differentiates every mode, the map
\(h(a)=v_{(-1)}a\) satisfies
\[
 dh(a)+hd(a)=(dv)_{(-1)}a=a.
\]
An exact vacuum therefore contracts the state complex.
In particular, a complex with \(H^0=0\) and nonzero
total cohomology cannot carry such a vacuum.

A geometric dual must retain both kinds of data seen
in the examples: simultaneous pole bounds, and
coefficients on partial collisions such as the
matrix context \(PQP\) and the scalar classes of
Theorem~\ref{e1f:thm:support}. For \(V\otimes B\),
charged Heisenberg families and the scalar vacuum,
Chapter~\ref{cl:chapter} constructs the regular
word-image system and the independent-input lattice,
with their consecutive collision maps. Field
evaluation relates these two objects and can have
a nonzero kernel. A full comparison must also
specify transitions between word orders and their
compatibility with field operations. Chapter~\ref{cqe:chapter}
states the additional coproduct and module-exchange
data on those operations.
